% Revised proposal built on the user's BIThesis 3.7.4 class, not the old thesis text.
% Compile with XeLaTeX, without shell escape. Generated bibliography uses the
% project's independently maintained 34-reference review rather than the old .bib.
\documentclass[type=master,twoside=false,ctex={fontset=fandol}]{bithesis}
\BITSetup{
 info={title={基于多维时序数据融合的短期证券市场风险预测},author={邱琨原},studentId={待填写},
 school={待填写},major={待填写},supervisor={待填写},degree={专业学位硕士},
 classification={待确认},UDC={待确认},defenseDate={待填写},
 keywords={实际损失,多维时序,公告抽取,机构持股,概率混合}},
 style={head={北京理工大学专业学位硕士开题报告},betterTimesNewRoman=true,
 pageVerticalAlign=top},misc={tabularFontSize=5}}
\IfFontExistsTF{Noto Serif CJK SC}{
 \setCJKmainfont{Noto Serif CJK SC}
 \setCJKsansfont{Noto Sans CJK SC}
 \setCJKmonofont{Noto Sans Mono CJK SC}
 \setCJKfamilyfont{zhsong}{Noto Serif CJK SC}
 \setCJKfamilyfont{zhhei}{Noto Sans CJK SC}
 \setCJKfamilyfont{zhkai}{Noto Serif CJK SC}
 \setCJKfamilyfont{zhfs}{Noto Serif CJK SC}
}{}
\usepackage{longtable,tabularx,makecell,float,placeins}
\usepackage{pgfplots,pgfplotstable}
\pgfplotsset{compat=1.18}
\usepgfplotslibrary{fillbetween}
\usetikzlibrary{arrows.meta,positioning,calc,fit,backgrounds,patterns,matrix}
% Original vector diagrams and historical aggregate plots; no old thesis assets.
% Palette validator: categorical CVD PASS; green contrast WARN relieved by
% black direct labels, legends, and border/line style distinctions.
\definecolor{seriesblue}{HTML}{2A78D6}
\definecolor{seriesorange}{HTML}{EB6834}
\definecolor{seriesgreen}{HTML}{1BAF7A}
\definecolor{inkgray}{HTML}{52514E}
\definecolor{rulegray}{HTML}{D8D8D8}
\tikzset{box/.style={draw=rulegray,rounded corners=2pt,fill=white,text=black,
 text width=2.9cm,minimum height=1.05cm,align=center,font=\sffamily\small,inner sep=5pt},
 tinybox/.style={box,text width=2.15cm},
 sb/.style={box,fill=seriesblue!7,draw=seriesblue},
 eb/.style={box,fill=seriesorange!7,draw=seriesorange},
 gb/.style={box,fill=seriesgreen!7,draw=seriesgreen,dashed},
 arr/.style={-{Stealth[length=2mm]},draw=inkgray,line width=.8pt},
 hint/.style={font=\sffamily\small,text=black,align=center}}
\pgfplotsset{reportaxis/.style={width=.94\linewidth,height=6cm,
 font=\sffamily\small,tick label style={font=\sffamily\footnotesize},
 label style={font=\sffamily\small},grid=major,grid style={rulegray,thin},
 axis line style={inkgray,thin},tick style={inkgray},scaled ticks=false,
 legend style={draw=none,font=\sffamily\small,text=black},
 every axis plot/.append style={line width=1.2pt}}}

\newcommand{\FigAsync}{
\begin{tikzpicture}[x=1.45cm,y=1.2cm]
 \foreach \i in {0,...,6}{\draw[rulegray](\i,.2)--(\i,3.25);\node[hint]at(\i,0){$t_{\i}$};}
 \node[hint,anchor=east]at(-.25,2.8){价格};
 \node[hint,anchor=east]at(-.25,1.7){公告};
 \node[hint,anchor=east]at(-.25,.6){持股};
 \foreach \i in {0,...,6}{\fill[seriesblue](\i,2.8)circle(2.8pt);}
 \fill[seriesorange](1,1.7)circle(3pt);\fill[seriesorange](4,1.7)circle(3pt);
 \draw[arr](1,1.7)--(2,1.7);\draw[arr](4,1.7)--(5,1.7);
 \draw[seriesgreen,thick](3,.6)rectangle++(.12,.12);\draw[arr](3,.6)--(4,.6);
 \node[hint,above]at(1,1.7){公开};\node[hint,above]at(4,1.7){新公告};
 \node[hint,below]at(3,.6){名单公开};
 \node[hint,text width=12cm]at(3,-.6){蓝圆点：价格更新；橙圆点：公告；绿方框：名单。\  箭头表示公开后的首个市场日可用，非实测频率。};
\end{tikzpicture}}

\newcommand{\FigAvailability}{
\begin{tikzpicture}[node distance=9mm and 8mm]
 \node[tinybox](r){报告期 $R$\\描述的业务期间};
 \node[tinybox,right=of r](p){公开日 $P$\\首次可获取日期};
 \node[tinybox,right=of p](v){可用日 $v$\\$P$ 后首个市场日};
 \node[tinybox,right=of v](t){预测日 $t$\\查询 $v\le t$};
 \draw[arr](r)--(p);\draw[arr](p)--(v);\draw[arr](v)--(t);
 \node[box,below=12mm of p,text width=7.7cm](version){更正公开形成新版本；原版保留独立键\\行业等资料还需同时满足其业务生效日};
 \draw[arr](p)--(version);\draw[arr](version.east)-|(t.south);
\end{tikzpicture}}

\newcommand{\FigArchitecture}{
\begin{tikzpicture}[node distance=7mm and 9mm]
 \node[sb](s0){个股GRU64\\市场GRU32、行业GRU32};
 \node[sb,right=of s0](s1){128维共同价格上下文\\历史值与有效掩码};
 \node[sb,right=of s1](s2){价格专家 $\pi_s$\\四档softmax};
 \node[eb,below=of s0](e0){冻结Qwen正文\\预测字段与原生字段};
 \node[eb,right=of e0](e1){价格上下文query\\事件选择与聚合};
 \node[eb,right=of e1](e2){事件专家 $\pi_e$\\接收相同价格上下文};
 \node[gb,below=of e0](g0){接受的公开完整名单\\比例、边龄与历史身份};
 \node[gb,right=of g0](g1){邻居价格+事件状态\\机构消息去目标自身};
 \node[gb,right=of g1](g2){关系专家 $\pi_g$\\接收相同价格上下文};
 \foreach \a/\b in {s0/s1,s1/s2,e0/e1,e1/e2,g0/g1,g1/g2}{\draw[arr](\a)--(\b);}
 \draw[arr](s1.south)--(e1.north);
 \node[box,below=11mm of g1,text width=9.7cm](out){可靠性权重 $\alpha_s,\alpha_e,\alpha_g$ 共用于四档\\$\pi=\alpha_s\pi_s+\alpha_e\pi_e+\alpha_g\pi_g$\\$p_{loss},p_5,p_{10}$；全缺模态回价格专家};
 \draw[arr](s2.east)--++(.35,0)|-(out.east);
 \draw[arr](e2.east)--++(.2,0)|-(out.east);
 \draw[arr](g2.south)|-(out.east);
 \node[hint,below=5mm of out]{图例：蓝实线框为价格；橙实线框为事件；绿虚线框为关系。};
\end{tikzpicture}}

\newcommand{\FigProducts}{
\begin{tikzpicture}[node distance=9mm and 8mm]
 \node[box](index){时点上市总体 $U_t$\\证券、日期与当时状态};
 \node[sb,below left=of index](raw){市场日序列索引\\三类价格与掩码};
 \node[eb,below=of index](events){事件版本缓存\\原文哈希与可用日};
 \node[gb,below right=of index](edges){完整名单版本流\\主体、比例与边龄};
 \node[box,below=16mm of events,text width=8.7cm](batch){日装配：目标与邻居去重，只读取过去信息\\离线结果左连接；已知标签与未知原因另返回};
 \draw[arr](index)--(raw);\draw[arr](index)--(events);\draw[arr](index)--(edges);
 \draw[arr](raw.south)|-(batch.west);\draw[arr](events)--(batch);\draw[arr](edges.south)|-(batch.east);
\end{tikzpicture}}

\newcommand{\FigLabel}{
\begin{tikzpicture}[node distance=9mm and 9mm]
 \node[box](u){预测日上市总体 $U_t$\\由当时身份状态决定};
 \node[box,right=of u](calendar){全市场日历\\$t+1,\ldots,t+5$};
 \node[box,right=of calendar](r){未来认证收益\\$R^{(5)}=\prod_k(1+r_k)-1$};
 \node[box,below=of u](all){每个成员都保留预测\\短历史使用掩码};
 \node[box,below=of calendar](join){结果左连接原预测索引\\不可观测则记录原因};
 \node[box,below=of r](bins){四档：非亏损、普通亏损\\较大亏损、严重亏损};
 \draw[arr](u)--(calendar);\draw[arr](calendar)--(r);\draw[arr](u)--(all);
 \draw[arr](all)--(join);\draw[arr](r)--(join);\draw[arr](join)--(bins);
 \node[hint,below=9mm of join,text width=10cm]{离线结果不决定 $U_t$ 成员；不减市场/行业/因子\\不要求252日回归预热；退市结果未知不填0或$-100\%$。};
\end{tikzpicture}}

\newcommand{\FigNamespace}{
\begin{tikzpicture}[node distance=9mm and 10mm]
 \node[box](cc){CompanyCode\\原公司编码空间};
 \node[box,right=of cc](lc){ComCd\\上市公司编码空间};
 \node[box,right=of lc](stock){A股证券身份\\历史有效区间};
 \node[box,below=of lc,text width=5cm](bridge){跨表映射与来源凭证\\观察起止不等于上市/退市日期};
 \node[box,below=of cc](holder){GDID：股东/产品\\不是目标证券};
 \node[box,below=of stock](fund){SecuCd：股东所属证券\\不是被持股A股代码};
 \draw[arr](cc)--node[above,hint]{分别保留}(lc);
 \draw[arr](lc)--node[above,hint]{时点核准}(stock);
 \draw[arr](cc.south)|-(bridge.west);\draw[arr](stock.south)|-(bridge.east);
\end{tikzpicture}}

\newcommand{\FigAnnotation}{
\begin{tikzpicture}[node distance=8mm and 11mm]
 \node[box](sample){锁定1500条、时间/家族隔离\\训练800 / 验证200 / 测试500};
 \node[box,below left=of sample](a){A独立试标200\\训练子集};
 \node[box,below right=of sample](b){B独立试标42\\既定双标交集};
 \node[box,below=19mm of sample](freeze){比较独立原稿，讨论后\\只据训练试标冻结规范};
 \node[box,below left=of freeze](resta){A继续1300\\主样本不重新抽取};
 \node[box,below right=of freeze](restb){B继续258\\合计双标300};
 \node[box,below=19mm of freeze,text width=7.7cm](gold){裁决绑定独立原稿哈希与人工放行\\SFT训练与独立语义评价分别导出};
 \draw[arr](sample)--(a);\draw[arr](sample)--(b);\draw[arr](a.south)|-(freeze.west);\draw[arr](b.south)|-(freeze.east);
 \draw[arr](freeze)--(resta);\draw[arr](freeze)--(restb);\draw[arr](resta.south)|-(gold.west);\draw[arr](restb.south)|-(gold.east);
\end{tikzpicture}}

\newcommand{\FigPooling}{
\begin{tikzpicture}[node distance=7mm and 5mm]
 \node[tinybox](system){系统/任务提示\\排除};
 \node[tinybox,right=of system](wrap){公司包装/标记\\排除};
 \node[sb,right=of wrap](body){正文token\\位置掩码池化};
 \node[tinybox,right=of body](answer){生成JSON/padding\\排除};
 \node[box,below=12mm of body,text width=7.8cm](vector){冻结正文向量与预测结构分别保存\\JSON失败不自动抹去合法正文表示};
 \draw[arr](body)--(vector);
\end{tikzpicture}}

\newcommand{\FigCache}{
\begin{tikzpicture}[node distance=8mm and 9mm]
 \node[box](text){原文哈希\\家族与公开版本};
 \node[box,right=of text](model){基座/adapter revision\\分词器与模板};
 \node[box,right=of model](policy){正文位置/长度策略\\schema与解析版本};
 \node[box,below=12mm of model,text width=7.8cm](key){组合缓存键 $K_e$\\正文、预测字段、各质量状态与失败原因};
 \node[box,below=of key,text width=7.8cm](query){日查询只取已可用内容\\任一组成改变重新生成相应缓存};
 \draw[arr](text.south)|-(key.west);\draw[arr](model)--(key);\draw[arr](policy.south)|-(key.east);\draw[arr](key)--(query);
\end{tikzpicture}}

\newcommand{\FigCertification}{
\begin{tikzpicture}[node distance=8mm and 9mm]
 \node[box](a){总股东1/1候选\\先区分名单口径};
 \node[box,right=of a](b){完整性、并列排名\\原版与更正对应};
 \node[box,right=of b](c){主体、证券与历史区间\\区分产品/管理人/托管};
 \node[gb,below=of c](d){接受的完整公开版本\\总股本比例入边};
 \node[box,below=of a](u){未决或冲突记录\\保存来源，隔离待核};
 \draw[arr](a)--(b);\draw[arr](b)--(c);\draw[arr](c)--(d);
 \draw[arr](b.south)|-(u.east);\draw[arr](a)--(u);
 \node[hint,below=9mm of d,text width=9.7cm]{图例：绿虚线框为接受的关系输入。\\辅助记录未接受不删除预测成员，关系掩码置0。};
\end{tikzpicture}}

\newcommand{\FigBipartite}{
\begin{tikzpicture}[node distance=10mm and 14mm]
 \node[box](i){目标 $i$\\状态 $x_i=[c_i;e_i;M_{e,i}]$};
 \node[box,below=of i](k){共享邻居 $k$\\已知价格+事件状态 $x_k$};
 \node[box,below=of k](l){共享邻居 $\ell$\\已知价格+事件状态 $x_\ell$};
 \node[box,right=of k](j){机构 $j$\\$S_j=\sum_k w_{kj}\phi(x_k)$};
 \node[gb,right=of j,text width=3.9cm](u){送回目标的消息\\$\displaystyle u_{j\setminus i}=\frac{S_j-w_{ij}\phi(x_i)}{d_j-w_{ij}}$\\无其他正权邻居则不可用};
 \draw[arr](i.east)-|(j.north);\draw[arr](k)--(j);\draw[arr](l.east)-|(j.south);\draw[arr](j)--(u);
 \node[hint,below=8mm of u,text width=8cm]{绿虚线框：关系消息；只传播过去状态，不传播未来标签。};
\end{tikzpicture}}

\newcommand{\FigEnvelope}{
\begin{tikzpicture}
\begin{axis}[reportaxis,height=6.7cm,xmin=0,xmax=2,ymin=0,ymax=1.12,
 xlabel={归一化披露年龄 $A$},ylabel={原始概率混合的允许权重},
 xtick={0,.5,1,1.5,2},ytick={0,.25,.5,.75,1},legend pos=north east]
 \addplot[seriesblue,domain=0:2,samples=100]{exp(-x)};
 \addlegendentry{单分支上界 $u\le e^{-A}$}
 \addplot[seriesorange,dashed,domain=0:2,samples=100]{2*exp(-x)-exp(-2*x)};
 \addlegendentry{双分支联合上界 $g$}
 \node[font=\sffamily\footnotesize,text=black,anchor=west]at(axis cs:1.3,.29){单分支};
 \node[font=\sffamily\footnotesize,text=black,anchor=west]at(axis cs:1.3,.55){双分支};
\end{axis}\end{tikzpicture}}

\newcommand{\FigTraining}{
\begin{tikzpicture}[node distance=8mm and 9mm]
 \node[box](read){当日目标与邻居并集\\去重读取过去资料};
 \node[box,right=of read](encode){当前GRU/投影参数\\共享市场与行业编码};
 \node[box,right=of encode](message){上下文事件注意力\\价格+事件去自身消息};
 \node[box,below=of message](fusion){可靠性混合四档概率\\匹配的模态dropout};
 \node[box,left=of fusion](loss){已认证训练目标\\混合NLL与共同正则};
 \node[box,left=of loss](update){反向传播和更新\\不监督邻居未来标签};
 \draw[arr](read)--(encode);\draw[arr](encode)--(message);\draw[arr](message)--(fusion);\draw[arr](fusion)--(loss);\draw[arr](loss)--(update);
 \draw[arr](update.north)|-(encode.south);
\end{tikzpicture}}

\newcommand{\FigFolds}{
\begin{tikzpicture}[x=1.4cm,y=1.08cm,font=\sffamily\small]
 \foreach \yr [count=\c from 0] in {2017,2018,2019,2020,2021,2022,2023,2024}{\node at(\c+.5,4.3){\yr};}
 \foreach \r in {0,1,2,3}{
   \pgfmathtruncatemacro{\endtrain}{2+\r}
   \foreach \c in {0,...,\endtrain}{\fill[black!7](\c,3-\r)rectangle(\c+.95,3.8-\r);}
   \pgfmathsetmacro{\val}{3+\r}\pgfmathsetmacro{\test}{4+\r}
   \fill[black!18](\val,3-\r)rectangle(\val+.62,3.8-\r);
   \fill[black!32](\val+.65,3-\r)rectangle(\val+.95,3.8-\r);
   \draw[inkgray,thick](\test,3-\r)rectangle(\test+.95,3.8-\r);
   \node[font=\sffamily\scriptsize]at(\val+.31,3.4-\r){选参};
   \node[font=\sffamily\scriptsize,rotate=90]at(\val+.8,3.4-\r){校准};
   \node at(\test+.47,3.4-\r){测试};\node at(1.1,3.4-\r){训练};
 }
 \node[align=left,anchor=north west,text width=12.5cm]at(0,-.4){图例：浅灰训练；中灰验证1--8月选参；深灰9--12月校准；边框测试。\\2024测试至7月24日；所有交界按实际结果结束日purge。};
\end{tikzpicture}}

\newcommand{\FigEvaluation}{
\begin{tikzpicture}[node distance=9mm and 9mm]
 \node[box](rows){同总体、同可评价成员\\模型与对照逐行配对};
 \node[box,right=of rows](seed){每折每种子先算指标\\再取五种子均值};
 \node[box,right=of seed](equal){四折等权差值\\年度和覆盖并列};
 \node[box,below=of seed](blocks){重叠的移动日期块\\整日公司与双方一起抽};
 \node[box,below=of equal](ci){20日主块、5/60敏感性\\1000次配对区间};
 \draw[arr](rows)--(seed);\draw[arr](seed)--(equal);\draw[arr](rows.south)|-(blocks.west);\draw[arr](blocks)--(ci);\draw[arr](equal)--(ci);
\end{tikzpicture}}

\newcommand{\FigDependencies}{
\begin{tikzpicture}[node distance=10mm and 9mm]
 \node[box](base){历史证券/状态与收益口径\\时点总体和实际损失候选};
 \node[box,below left=of base](human){锁定人工金标\\独立语义评价};
 \node[box,below right=of base](graph){名单、主体、公开版本\\历史证券对应};
 \node[box,below=18mm of base](cpu){真实取窗、未知原因、日装配与强对照\\CPU边界测试、校准与恢复接口};
 \node[box,below=of cpu](gpu){认证条件满足后GPU抽取/缓存\\四折五种子正式比较与统计};
 \draw[arr](base)--(human);\draw[arr](base)--(graph);\draw[arr](human.south)|-(cpu.west);\draw[arr](graph.south)|-(cpu.east);\draw[arr](cpu)--(gpu);
\end{tikzpicture}}

\newcommand{\FigOutput}{
\begin{tikzpicture}[node distance=9mm and 10mm]
 \node[box](input){当日上市总体与有效资料\\来源、身份和状态};
 \node[box,right=of input](model){最终四档校准概率\\$p_{loss},p_5,p_{10}$};
 \node[box,right=of model](alert){绝对阈值预警\\Top10\%仅次要名单};
 \node[box,below=of model,text width=8.8cm](audit){市场/行业平均$p_5,p_{10}$：预期损失成员占比\\专家概率、混合权重、版本与未知结果原因};
 \draw[arr](input)--(model);\draw[arr](model)--(alert);\draw[arr](model)--(audit);
\end{tikzpicture}}

\newcommand{\FigMarketIndustry}{
\begin{tikzpicture}[node distance=10mm and 9mm]
 \node[sb](security){个股：实际收益\\过去累计/波动/下行};
 \node[sb,right=of security](market){市场：认证原始收益\\FF3与过去状态};
 \node[sb,right=of market](industry){行业：时点已知门类\\行业收益与过去状态};
 \node[box,below=of security](si){掩码GRU64\\短历史仍保留};
 \node[box,below=of market](mc){掩码GRU32\\每日期只编码一次};
 \node[box,below=of industry](ic){掩码GRU32\\每行业日期编码一次};
 \node[box,below=12mm of mc,text width=9.5cm](out){共同价格上下文 $c_i\in\mathbb R^{128}$\\进入价格/事件/关系全部专家与事件query\\未来实际收益标签保留共同下跌，不做因素扣除};
 \draw[arr](security)--(si);\draw[arr](market)--(mc);\draw[arr](industry)--(ic);
 \draw[arr](si.south)|-(out.west);\draw[arr](mc)--(out);\draw[arr](ic.south)|-(out.east);
\end{tikzpicture}}

\newcommand{\FigYearCounts}{
\begin{tikzpicture}\begin{axis}[reportaxis,ybar,bar width=7pt,
 xlabel={历史审计源记录的公开年份},ylabel={股东记录数（万条）},ymin=0,ymax=48,
 symbolic x coords={2015,2016,2017,2018,2019,2020,2021,2022,2023,2024,2025,2026},
 xtick=data,x tick label style={rotate=45,anchor=east,font=\sffamily\footnotesize},
 ytick={0,10,20,30,40},enlarge x limits=.06]
 \addplot[fill=seriesblue,draw=none]table[x=year,y=wan,col sep=space]{data/yearly.dat};
 \node[font=\sffamily\footnotesize,text=black,anchor=south]at(axis cs:2023,43.1){430,988};
 \node[font=\sffamily\footnotesize,text=black,anchor=south]at(axis cs:2026,.04){340};
\end{axis}\end{tikzpicture}}

\newcommand{\FigMissing}{
\begin{tikzpicture}\begin{axis}[reportaxis,xbar,bar width=9pt,
 xmin=0,xmax=66,xlabel={占历史审计全量股东记录的比例（\%）},
 symbolic y coords={信息来源,A股代码,股东ID,流通股比例},ytick=data,
 yticklabel style={font=\sffamily\small},xtick={0,20,40,60},enlarge y limits=.25,
 nodes near coords,point meta=explicit symbolic,every node near coord/.append style={font=\sffamily\footnotesize,text=black,anchor=west}]
 \addplot[fill=seriesblue,draw=none]table[x=pct,y=field,meta=label,col sep=space]{data/missing.dat};
\end{axis}\end{tikzpicture}}

\newcommand{\FigUnknownBounds}{
\begin{tikzpicture}[x=1cm,y=1cm]
 \draw[inkgray,thin](0,0)--(11,0);
 \foreach \x/\lab in {0/0,3/{B/N},8/{(B+M)/N},11/1}{\draw[inkgray](\x,-.12)--(\x,.12);\node[hint,below]at(\x,-.15){$\lab$};}
 \draw[seriesblue,line width=2pt](3,0)--(8,0);
 \node[hint,above,text width=7.5cm]at(5.5,.5){未知成员可能均非正例，也可能均为正例\\实际风险广度在此区间内};
 \node[box,text width=10cm]at(5.5,-1.7){$N$：预测总体人数；$B$：已知较大/严重亏损人数\\$M$：结果未知人数；区间宽度 $M/N$ 反映覆盖缺口};
\end{tikzpicture}}

\newcommand{\FigLogitCap}{
\begin{tikzpicture}\begin{axis}[reportaxis,ybar,bar width=12pt,
 ymin=0,ymax=.105,ytick={0,.025,.05,.075,.1},ylabel={固定价格主干时的最大概率},
 symbolic x coords={价格主干,单支路上限,双支路上限},xtick=data,
 xlabel={旧单位logit修正的解析边界}]
 \addplot[fill=seriesblue,draw=none]coordinates{(价格主干,.01)(单支路上限,.02672363)(双支路上限,.06945316)};
 \node[hint,above]at(axis cs:价格主干,.01){0.0100};
 \node[hint,above]at(axis cs:单支路上限,.02672363){0.0267};
 \node[hint,above]at(axis cs:双支路上限,.06945316){0.0695};
\end{axis}\end{tikzpicture}}

\newcolumntype{Y}{>{\raggedright\arraybackslash}X}
\newcommand{\method}{可靠性约束的四档概率混合}
\newcommand{\source}[1]{\par\smallskip{\small\sffamily\color{inkgray}来源与口径：#1}\par}
\newcommand{\reporttable}[3]{\begin{table}[tbp]\centering\caption{#1}\label{#2}\small #3\end{table}}
\newcommand{\reportfigure}[4]{\begin{figure}[tbp]\centering\begingroup\setbox0=\hbox{#3}\ifdim\wd0>\linewidth\resizebox{\linewidth}{!}{\box0}\else\box0\fi\endgroup\caption{#1}\label{#2}\source{#4}\end{figure}}
\setlength{\parindent}{2em}
\setlength{\parskip}{5.5pt}
\setlength{\textfloatsep}{16pt plus 2pt minus 2pt}
\setlength{\intextsep}{14pt plus 2pt minus 2pt}
\captionsetup{font=small,labelfont=bf}
\AtBeginDocument{\renewcommand{\normalsize}{\fontsize{12bp}{22bp}\selectfont}\setstretch{1}\normalsize}
\makeatletter
\newcommand{\ReportTypesettingCheck}{\typeout{REPORT_BODY_FONT_SIZE_PT=\f@size}\typeout{REPORT_BODY_BASELINE_PT=\the\baselineskip}}
\makeatother
\hypersetup{pdftitle={基于多维时序数据融合的短期证券市场风险预测——专业学位硕士开题报告},
 pdfauthor={邱琨原},pdfsubject={两年制专业学位硕士；研究设计与已核数据分析},hidelinks}
\begin{document}
\begin{titlepage}
 \centering
 \vspace*{8mm}
 {\heiti\zihao{1}北京理工大学\par}
 \vspace{10mm}
 {\heiti\zihao{1}专业学位硕士\par}
 \vspace{4mm}
 {\heiti\zihao{1}学位论文开题报告\par}
 \vspace{12mm}
 {\heiti\zihao{2}基于多维时序数据融合的\par 短期证券市场风险预测\par}
 \vspace{12mm}
 {\zihao{4}\renewcommand{\arraystretch}{1.25}
 \begin{tabular}{rl}
 姓名：&邱琨原\\ 学制：&2年\\ 入学年级：&待填写\\ 学号：&待填写\\
 类别、领域：&待填写\\ 培养单位：&待填写\\
 校内导师：&待填写\\ 拟申请学位方式：&学位论文\\
 正式开题日期：&待填写
 \end{tabular}\par}
 \vfill
 {\zihao{4}材料更新：2026年11月26日\par}
\end{titlepage}
\frontmatter
\tableofcontents
\clearpage
\mainmatter
% BIThesis mainmatter sets linespread=1.53 for its default 14.4pt baseline.
% Our explicit 22bp baseline must not be multiplied again.
\setstretch{1}
\fontsize{12bp}{22bp}\selectfont
\ReportTypesettingCheck
\chapter{选题依据}
\section{实践背景：风险审核面对的是不断变化的信息集}
上市公司的短期风险筛查通常从一个很具体的问题开始：在当天可获取的资料中，哪些公司值得优先复核？审核人员可以看到近期收益、经营预告、诉讼进展和持股名单，却难以每天逐项阅读所有公司的材料。市场行情不断更新，公告出现的时间不规则，持股名单又带有明显的报告期和披露滞后。把这些资料整理为公司日风险排序，首先需要解决的并不是网络层数，而是资料在历史时点上的可用性。

业务日期和信息日期的区别，在低频披露中尤其重要。季度末描述的是持股状态对应的期间，公开日描述的是市场参与者何时能够获知该状态，数据库更新时间则可能只是更正或补录发生的时间。三者混用会产生两类问题。一类是把晚公布的内容提前放入回测；另一类是资料虽然确实已公开，却在很长时间内被当作同样新鲜的信息使用。前者影响实验有效性，后者影响模型对辅助信息的信任程度。图\ref{fig:async}把两类问题放在同一时间框架下。

\reportfigure{收益、公告与持股关系的异步更新}{fig:async}{\FigAsync}{研究流程示意。时间点仅表达先后次序，不是实际公司个案，也不表示经统计得到的更新频率。}

仅靠日期过滤不能解决全部问题。某份名单在预测日前已公开，仍可能因公司经营变化、机构调仓或名单更正而逐渐失去代表性。文本抽取也存在类似情形：模型可以输出一个格式正确的答案，答案却未必准确概括经营变化。若融合模型允许这类资料大幅改变主干分数，新增信息可能不但没有帮助，还会增加排序波动。因此，资料的存在、可追溯程度和陈旧程度需要进入模型设计，而不是只保存在数据清洗日志里。

本课题面向上述审核场景，研究A股公司未来五个市场交易日的条件左尾事件。主干使用持续更新的收益和三因子序列；业绩预告语义、共同机构持股关系作为异步辅助资料；诉讼信息形成可直接追溯的结构控制。研究输出是离线风险排序方法和原型，供有限审核容量下的风险筛查研究使用。模型不判断某条公告是否因果性地引起下跌，也不把共同持股解释为已经识别的风险传染机制。

\section{研究对象与任务边界}
“尾部”需要在任务中给出可以重复计算的定义。普通涨跌分类把零收益作为分界，而年度崩盘研究常使用较长时期的偏度或崩盘代理变量。本课题选择未来五日累计因子残差相对公司自身历史分布的左尾。这样能够将近期价格状态和公司特有异常波动联系起来，也允许不同公司使用适应自身历史分布的阈值。它不是监管定义的风险等级，更不是破产概率。

预测时点定为交易日收盘。收益、因子及公告关系的查询都以该时点为界，未来收益只用于离线定义结果标签。没有可靠时分秒的披露资料，从公开后的下一市场交易日起进入特征。采用这一约定会损失部分当天已公开的信息，但可以避免把盘后材料当作盘中已知事实。对于目前以日期为粒度的源数据，这是一项明确、保守且可以程序验证的取舍。

图\ref{fig:availability}还区分了更正内容。更正不是把原始记录的一个数值“修好”以后再覆盖所有历史日期，而是新的可见版本。训练与测试的历史查询必须返回当时能够使用的内容。若源表无法提供旧版本，原公告比对和数据隔离就成为必要环节。资料认证的结果会影响辅助模态覆盖，但不应改变同一个公司日标签是否参加模型比较。

\reportfigure{报告期、公开日、生效日和更正版本的关系}{fig:availability}{\FigAvailability}{按日期级披露拟定的可用性规则。更正保留独立版本，不据报告期回填。}

\section{相关研究为本课题提供了什么}
\subsection{崩盘与公司特有收益}
Chen、Hong和Stein从条件偏度讨论股票崩盘预测，强调过去市场状态与极端下跌分布之间的联系\cite{r1}。Jin和Myers关注信息不透明、收益共同波动及风险积累\cite{r2}，Hutton等将财务报告不透明性与崩盘风险联系起来\cite{r3}。这些研究使公司特有信息成为尾部研究的重要线索，但其年度变量、经济解释和样本结构与本课题的日频预测不同。

Fama和French三因子模型提供了一个熟悉且可计算的共同风险参照\cite{r4}。这里使用三因子并非认为它已经解释全部收益变化，而是为标签建立一致口径：预测起点估计系数，冻结系数后观察未来残差，再与过去已结束窗口比较。这一操作定义有利于复核，但残差仍会包含遗漏因子、估计误差和价格噪声。实验结论应始终对应这一标签，不把它扩写为“识别所有形式的金融风险”。

\subsection{文本从情绪走向事件与证据}
Tetlock及其后续研究说明新闻语言可能包含数值变量之外的信息\cite{r5,r6}。Loughran和McDonald指出财务语境会改变词语的含义\cite{r7}。这些线索对业绩预告尤其直接：盈利但同比下降、仍亏损但亏损收窄，分别涉及绝对状态和变化方向，不能只根据“盈利”或“亏损”两个词判断。

EFSA将中文财经新闻分析细化到公司、行业、事件和情感的组合\cite{r24}。本课题借鉴事件级分析思路，但把正式预告的抽取目标限定为方向、明确原因和原文证据。选择有限字段，一方面使人工规范能够执行，另一方面便于追查模型为什么产生某个事件表示。预训练模型、LoRA与QLoRA提供的是表示和适配工具\cite{r14,r15,r16,r21,r22,r23}；是否适合本地公告，还需要独立金标评价。

\subsection{关系网络与融合模型的直接近邻}
Feng等研究股票关系下的时序排序\cite{r8}，Cheng等研究多模态图神经网络用于金融序列\cite{r9}，Hou和Liu直接讨论机构交叉持股与崩盘风险\cite{r10}。关系与风险的联系因此不是空白。GraphSAGE、图卷积和图注意力提供了多种邻居聚合方式\cite{r17,r18,r19}，但这些算子并不自动解决持股名单来源和金融主体身份的问题。

Jing等把改造Graph WaveNet用于中国股票崩盘风险\cite{r11,r20}；MSGCA同时利用指标、动态文本和图关系，并设置门控与主干回接\cite{r12}。两者要求本课题分别面对任务方向和融合架构上的近邻，而不能仅与一个弱线性分类器比较。表\ref{tab:neighbors}将这种关系具体化。不同论文的频率、标签和信息集并不一致，原论文分数不能直接排入本地结果表。

\reporttable{主要近邻与本任务需要重新验证的部分}{tab:neighbors}{
\begin{tabularx}{\linewidth}{p{2.65cm}YY}\toprule
研究线索&已有内容&本课题中的比较问题\\\midrule
经典崩盘研究&条件偏度、信息透明度、公司治理&年度解释性变量能否转化为只用历史的日频特征\\
文本与EFSA&财经语言和事件级语义&正式预告三字段是否抽取准确、是否有下游增量\\
共同持股&机构交叉持股与崩盘风险联系&公开时点图是否超过收益与公告信息\\
Graph WaveNet&图时序建模及中国股票崩盘应用&同一五日标签、同一时间折的任务适配表现\\
MSGCA&指标、文本、图的门控融合&同编码器下年龄幅度约束是否仍有必要\\\bottomrule
\end{tabularx}}

\section{本课题的研究切入点}
相近文献已经拥有三路输入、门控、关系建模和收益主干。本课题的切入点不是再增加一个组件，而是把辅助信息对主干分数的作用限制为可用性、质量、年龄和有界内容项的乘积。这种设计可以先给出修正幅度上界，再在相同信息条件下检验其代价和收益。

关键问题是约束是否过强。稳定的持股关系可能在较长时间里仍有效，一份旧公告也可能解释持续经营变化。年龄越大不等于信息必然无用，所以衰减率需要由训练学习，同时设置年龄仅作普通特征、去全部年龄条件和仅去边级衰减的控制。它们可以区分“增加年龄信息”与“限制修正幅度”，也能检验图内相对权重和图外绝对陈旧度分别起什么作用。

专业学位研究还需要把方法落实到可以检查的工程对象。公司编码、公开版本、正文缓存和日批次索引都应有清楚的接口与失效条件。其价值在于让风险排序能追溯到具体资料与模型版本，而非只保留一个分数。本文后续据此展开数据、模型、实验和实现设计；独立文献综述另行提交，不把本章当作对完整综述的重复粘贴。
\ReportTypesettingCheck\FloatBarrier
\chapter{研究内容}
\section{总体思路：以实际损失为共同监督目标}
研究围绕一个共同目标组织不同信息，而不是为价格、公告和图分别定义互不相同的风险。预测日$t$的证券$i$具有过去价格历史、市场与行业环境，以及已公开事件和持股关系。模型输出未来五个市场日累计收益的四档分布，由此得到发生亏损、达到$5\%$损失和达到$10\%$损失的嵌套概率。不同来源对同一未来结果提供条件信息，它们不承担彼此替代的标签。

价格信息既是持续更新的底座，也是理解公告和关系的环境。同一经营预告在上涨、普跌或高波动时期，可能对应不同后续损失概率；共同机构下的多只证券，也可能主要受到同一行业冲击。本研究将个股、市场和行业分别编码成64、32和32维历史状态，拼接得到128维共同价格上下文。事件注意力以该上下文为query，关系消息读取邻居截至当日已知的价格与事件状态，三个专家均接收共同环境。

\reportfigure{共同价格上下文与可靠性约束的四档专家混合}{fig:architecture}{\FigArchitecture}{批准研究设计。各维度、掩码和概率接口为拟实施合同，图不表示正式风险模型已经训练；蓝、橙、绿框均有直接文字及边框编码。}

图\ref{fig:architecture}体现了两个设计选择。其一，价格专家不只是输出一个等待小幅修补的logit，而是一个完整的四档专家；事件和关系专家也能表达强损失情形。其二，辅助专家的最大权重由资料可靠性限制，而不是由它们自身给出的风险高低决定。强事件能否主导，取决于可用、质量、年龄和学习门，不由任意固定logit帽限制。这一结构仍属于混合专家思路，创新假设集中在时点资料约束的具体构造和检验。

\section{研究问题与证据安排}
首先研究共同价格状态对实际损失的预测能力。个股历史波动与共同市场环境都有合理经济依据，但GRU不必然超过滚动收益和下行统计。研究将保留透明Logit和强树，比较个股价格、价格加市场行业以及完整价格上下文模型，分析共同因素是否尤其有助于普跌日期。去市场或去行业控制同时关闭相应显式派生特征，不允许删除一个分支后继续访问同一信息。

其次研究公告语义和持股关系的增量。公告方向是经营语义，不是未来收益标签；持股关系是公开的信息关联，不是已识别的因果传播。原生字段、预测结构和正文表示分层比较，可以区分直接数值、生成抽取与连续语义的作用。真实关系、去关系和保持度数及行业组成的重连，可以检查图增量是否只是网络容量或行业构成。各项比较使用相同预测总体和相同已知结果评价集合。

第三研究可靠性约束相对普通门控、拼接及标准概率MoE的净作用。约束可以给出缺失回退和年龄相关的原始概率偏移界，但这些性质不是泛化保证。正式比较需要观察AP、NLL、Brier、校准与预警漏报，并对自然年龄、低质量和低共享邻居分组分析。若约束只在一种覆盖群体有效，研究应解释条件而不是仅给跨年平均增幅。

\reporttable{研究问题与主要证据对象}{tab:questions}{
\begin{tabularx}{\linewidth}{p{2.8cm}YY}\toprule
研究问题&主要比较&需要观察的证据\\\midrule
共同价格环境的作用&透明统计、强树、个股与市场行业GRU&逐年$p_5$排序、概率质量、共同下跌漏报\\
公告语义的增量&原生字段、预测结构、正文、去事件&独立语义误差、覆盖与同信息下游差值\\
共享关系的增量&去关系、真实图、度数/行业组成保持重连&共享邻居、去自身边界与概率增量\\
可靠性约束的作用&拼接、普通门控、标准MoE与本方法&同目标校准、缺失回退、年龄界及压力\\
总体与结果分离&认证总体、观察候选及未知原因&全总体覆盖、监督子集与风险广度界限\\\bottomrule
\end{tabularx}}

\section{数据产品：预测索引与结果索引分离}
数据工作从原始资料的经济对象和历史时点开始。不同工作簿可能含同义表头，也可能使用不同公司码；原始名单可能混合总股东、流通股东及补充记录。标准化应保留源文件指纹、字段与行位置，用确定性映射形成候选产品。它的目的不是让所有缺失和冲突都变成“可用”，而是使合法连接与不确定状态可以明确识别。

预测索引以时点上市A股总体为主，记录证券、预测日、当时身份和交易状态。过去输入通过市场日序列索引、事件公开版本缓存和完整名单版本流按日查询。未来结果另表计算，左连接到预测索引后形成四档监督或未知原因。图\ref{fig:products}中的分离让模型输入不必通过一张已筛除未来缺失的宽表生成，也防止邻居节点隐含获得未来标签。

\reportfigure{时点总体、过去输入与离线结果的产品关系}{fig:products}{\FigProducts}{数据接口设计。预测成员由当时状态建立，结果可空；未来结果不参与历史输入查询和总体筛选。}

来源认证不足时，工程执行可以产生“观察总体/完整窗口候选面板”。这个名称反映能够从现有源行确认的范围，不声称覆盖所有当时上市证券。补齐股票主表、每日交易估值状态和终端价值之后，才判断哪些候选可以进入认证面板。候选数量需要实际运行生成，不能从原残差标签规模推算，更不能把原样本ID直接改名沿用。

\section{公告与关系表示的具体研究内容}
公告工作复用已锁定的1500条语料、800/200/500划分以及200试标和300双标。语义目标保持方向、明确原因和逐字证据，标注者不读取未来风险标签。Qwen完成适配和独立语义评价后冻结，风险任务使用其预测结构和正文向量，而非测试金标。正文池化排除系统、包装、padding和生成答案；缓存键将原文、模型、分词器、模板和正文范围绑定。

事件表示将合法正文与合法JSON分开。解析失败会影响结构字段，但未必使正文隐藏表示无效；证据可在正文找到，也不意味着它完整支持了语义答案。不同状态应进入质量记录和覆盖报告，不能合并成单一成功率。价格上下文query对最近可用事件进行选择，学习在当前环境下哪些事件更相关，最终事件专家仍预测同一四档结果。

关系工作以接受的完整公开名单为更新单位，区分主体与历史证券对应。新完整版本替换旧名单，Top10退出只说明离开披露范围，不等于清仓。比例和边龄形成非负权重，机构聚合邻居已知价格与事件状态，再扣除目标自身贡献。没有其他正权共享邻居时关系不可用。这样的接口才能支持“事件与价格状态沿关系聚合”的表述，只有传播价格向量的旧原型不应被描述为事件传播。

\section{训练、校准与输出之间的关系}
训练采用最终混合四档概率的NLL，并使用匹配的模态dropout和共同正则化。价格专家在辅助缺失样本上接受同一主目标的学习，事件与图专家也由混合目标联合优化。旧分支辅助BCE系数不进入主方案；若研究其影响，作为另行登记、同目标匹配的敏感性实验。这样，同编码器比较主要改变融合机制，而不是同时改变监督预算。

验证年1--8月用于选参和早停，9--12月用于最终四档温度校准及预警阈值，两个阶段按结果结束日剔除跨界窗口。温度作用于最终混合分布，不假定各专家先校准后混合自然可靠。全部档位使用同一组混合权重，校准后仍由四档求和获得嵌套概率。原始年龄界和非线性校准后的概率质量分别检查。

\reporttable{主要输入与输出的张量合同}{tab:tensors}{
\begin{tabularx}{\linewidth}{p{3.2cm}p{3.15cm}Y}\toprule
对象&形状或维度&含义\\\midrule
个股历史&$N\times60\times F_s$&实际收益及仅过去派生量，逐字段掩码\\
市场历史&$D\times60\times F_M$&市场原始收益、FF3及过去状态，每日共享\\
行业历史&$C\times60\times F_I$&时点已知门类环境，类别日期共享\\
共同价格上下文&$N\times128$&个股64、市场32、行业32拼接\\
正文与事件状态&$N\times K\times D_Q$&$K\le8$，正文与预测结构分别判定有效\\
二部边及邻居状态&$2\times E$；节点状态向量&比例、边龄、历史身份与价格+事件状态\\
专家与混合分布&三个$N\times4$及一个$N\times4$&非负归一，所有档位使用相同权重\\
结果与状态&$N$&四档监督可空，未知原因与监督掩码另存\\\bottomrule
\end{tabularx}}

表\ref{tab:tensors}的$F_s,F_M,F_I$由冻结字段顺序决定，不能沿用旧五字段序列或合成示例默认值来代表新装配。原始输入值、掩码、来源可用日和监督状态各自需要检查。最终原型保留全总体的$p_{loss},p_5,p_{10}$、专家分布和混合权重，按绝对阈值预警；每天Top10\%只是次要审核容量指标。市场和行业平均$p_5,p_{10}$用于估计损失成员占比，结合未知结果上下界评价共同下跌覆盖，不扩写为指数事件概率。
\ReportTypesettingCheck\FloatBarrier
\chapter{研究方案}
\section{数据连接与公开时点}
\subsection{研究范围及公司日主键}
日收益实值覆盖2015年1月5日至2024年7月31日，共9,091,641行。有效标签覆盖2017年2月6日至2024年7月24日。起始年份的差别来自历史估计和已结束阈值窗口的要求；结束日期的差别来自未来五市场日结果的需要。因子和公告表存在更晚记录，并不能据此延长风险测试。模型范围由收益结果可观测性决定，辅助资料再按公开时点查询。

日收益表中的证券代码和日期组成主键，当前该键已做唯一性核查。ComCd可以参与跨表连接，但不是证券交易代码；CompanyCode又属于另一个命名空间。旧转换器曾随工作簿列顺序选取两类公司码，造成编码混用。修正后，9,091,641行中4,318,444行仅公司码字段变化，收益、日期等标签依赖字段保持不变。这一发现说明来源解析错误和标签算法错误应分别检查，不能一发现公司码问题就宣布全部收益标签失效。

\reportfigure{公司、证券与股东所属证券的不同编码空间}{fig:namespace}{\FigNamespace}{字段连接设计；两命名空间独立保留。观察对应关系仅为候选，不是历史上市有效区间认证。}

图\ref{fig:namespace}中的股东所属证券编码还需要特别处理。基金或证券产品的SecuCd描述的是股东一侧，不是被持股公司的A股代码。若把它用来连接收益表，会产生结构上看似完整、实际指向错误对象的关系。历史证券对应应有有效日期和来源凭证，不能用最新全称或当前股票代码回填全部历史。

\subsection{日期级信息的查询规则}
设市场日历为有序集合$\mathcal T$。对只有公开日期$P_e$的事件，定义生效日
\begin{equation}
 v_e=\min\{t\in\mathcal T:t>P_e\}.
\end{equation}
在预测日$t$读取$v_e\le t$的版本，保存报告期作为业务解释字段，不将其当作生效时点。公告正文、原生数值和后续更正共同进入版本记录；后续更新不能修改首次日已经看到的内容。诉讼首次公开日适合描述事件开始，当前进展则按更新公开日进入控制变量。

公司最近事件年龄取市场日序号差，而不是自然日差。这样周末和长假不会在模型中变成额外收益步长。市场日历来自全市场日期，不来自该公司自己的非缺失收益日；停牌公司也沿同一日历查询事件。每个来源返回原始公开日、生效日、版本键和可用状态，装配器逐一验证最大可用日不晚于预测日。

\section{短期条件尾部标签}
\subsection{起点估计与未来结果}
令$r_{i,s}$为公司$i$日收益，$r_{f,s}$为无风险收益，$F_s$为流通市值加权市场溢酬、SMB和HML。先筛选交换/市场组合$2|A$，避免同一天不同因子组合重复进入连接。在$t$以前最近252个有限有效观测组成集合$\mathcal H_{i,t}$，估计
\begin{equation}
 r_{i,s}-r_{f,s}=\alpha_{i,t}+\beta_{i,t}^{\mathsf T}F_s+\epsilon_{i,s},\quad s\in\mathcal H_{i,t}.
\end{equation}
截距与三因子系数在定义该起点结果时冻结。未来结果为
\begin{equation}
 C_{i,t}^{(5)}=\sum_{k=1}^{5}\left[r_{i,t+k}-r_{f,t+k}-\widehat\alpha_{i,t}-\widehat\beta_{i,t}^{\mathsf T}F_{t+k}\right],
\end{equation}
其中$t+k$表示市场日历中的第$k$个后续交易日，不是第$k$个个股有效记录。未来五日任一必需观测缺失时，不产生该起点标签。

三因子调整使任务侧重公司相对共同市场风险的异常变化，但并不把残差解释为已识别的纯公司冲击。回归窗口、因子估计与遗漏变量都会影响结果。因此，实验使用同一标签比较模型，而不把年度崩盘研究或普通涨跌研究的指标直接用作精度参照。

\subsection{仅使用已结束历史窗口的阈值}
对每个历史起点$u$，其五日结果同样使用$u$时点冻结的系数。预测$t$的阈值只从结果结束日不晚于$t$的窗口选取最近252个有效值，记为$\mathcal C_{i,t}$。排序后固定取
\begin{equation}
 q_{.10,i,t}=C_{i,t,(26)},\quad q_{.05,i,t}=C_{i,t,(13)},\qquad
 Y_{i,t}^{(.10)}=\mathbf1\{C_{i,t}^{(5)}<q_{.10,i,t}\}.
\end{equation}
严格小于阈值才为正类。历史窗口不应以未来重新估计的系数回算；全样本分位数也不能代替起点历史分位数。图\ref{fig:label}呈现了系数和窗口的两条时间路径。

\reportfigure{起点系数、已结束窗口与五日结果的标签构造}{fig:label}{\FigLabel}{标签算法设计及已实施的计算口径；图中未来路径只属于离线监督，不进入预测输入。}

有效历史回归观测和有效结果窗口是两种要求。公司可能已经有252个可回归日，却还没有252个已结束五日窗口；某些窗口还会因停牌或缺失失效。阈值和结果都满足条件才形成共享目标样本。实际P10正类率不需要恰好等于10\%，因为未来分布会变化，窗口之间也有重叠。本课题在每折报告真实正类率，不用名义分位数替代样本统计。

\section{60市场日序列与训练期标准化}
序列主字段为日收益、超额收益和三个因子，得到$X_{i,t}\in\mathbb R^{60\times5}$。相同形状的$M_{i,t}$表示各字段是否有效。缺失项先填0，再与掩码拼接为十通道输入：
\begin{equation}
 \widetilde X_{i,t}=\big[X_{i,t}\odot M_{i,t};M_{i,t}\big],\quad
 h_{i,t}=\operatorname{GRU}_{64}(\widetilde X_{i,t}),\quad
 s_{i,t}=\tanh(W_sh_{i,t}+b_s).
\end{equation}
掩码让模型区分真实零收益与缺失占位。若向前找60个个股有效观测，长期停牌公司的窗口可能跨越远多于60个市场日，序列尺度将不再一致。表\ref{tab:windowcases}给出处理差异。

\reporttable{收益窗口中不同情况的处理}{tab:windowcases}{
\begin{tabularx}{\linewidth}{p{3cm}YY}\toprule
情况&历史输入&标签结果\\\midrule
有效零收益&值0，掩码1&按未来结果完整性计算\\
预测日收益缺失&缺失占位0，掩码0&其他条件满足时可保留目标\\
历史窗口部分缺失&保留相同市场日位置及字段掩码&回归历史按有效观测另行筛选\\
未来五日缺必需观测&未来不参与输入&不产生该起点标签\\
非市场交易日&不占收益序列位置&事件延到下一个市场日生效\\\bottomrule
\end{tabularx}}

各折标准化使用训练期有效值估计均值和尺度，验证与测试只应用已有参数。缺失填充值、异常截断界及正文PCA同样纳入版本记录。不能为方便把2017—2024全部数据一起拟合预处理。基础收益统计对照也遵守这一约定，避免深度模型和强树在不同信息条件下竞争。

GRU提供有序历史表示\cite{r13}，并不自动优于滚动统计。仅收益模型只使用其主BCE，不接收公告、图辅助监督或诉讼控制；同信息融合模型的主干则接收共同的诉讼结构变量。这一区分需要写进输入合同，不能只根据模型名称判断它实际看到了什么。

\section{预告金标：方向、原因和证据如何分开}
\subsection{语义目标与可执行规范}
标注对象是预告正文所表达的经营/财务变化。方向取改善、恶化、基本稳定、混合、不明确；原因是原文明确给出的自由短语数组；证据是1—3段逐字片段。没有原因时记空数组，不反过来把方向改成不明确。未填写、语义不明确和未提及原因分别属于进度状态、语义类别和原因状态。

表\ref{tab:annotationexamples}使用虚构短句说明这一区别。例句不代表真实公司，也不用于模型训练。对标注者而言，主要困难不是编写JSON，而是区分“仍亏损”与“变化改善”，以及判断一段话究竟陈述原因还是重复结果。原因字段不预设行业、市场或管理层分类，以免逼迫标注者补充正文没有给出的解释。

\reporttable{三字段标注的说明用例（全部虚构）}{tab:annotationexamples}{
\begin{tabularx}{\linewidth}{Yp{1.7cm}Y}\toprule
说明文本&方向&原因与证据处理\\\midrule
预计净利润同比增长，主要由于产品销量增加。&改善&原因：产品销量增加；证据复制原句\\
仍为亏损，但亏损额较上年减少。&改善&未提明确原因，原因数组为空；保留原句\\
预计盈利，但净利润比上年下降。&恶化&绝对盈利不能代替同比方向；原因按原文有无填写\\
净利润增长，但主营业务利润下降。&混合&重要指标明确反向，证据同时覆盖两者\\\bottomrule
\end{tabularx}}

\subsection{锁定样本、独立双标与裁决}
主样本1500条按2017—2019训练800、2020验证200、2021—2024年7月测试500固定。200试标是训练子集；300双标是主样本子集，二者交集42。先由A标200、B独立标42，在原稿上计算一致率及方向Kappa，再讨论和冻结规范。余下A1300与B258按同一版本继续，避免试标反馈来自验证/测试而改变规则。

\reportfigure{固定成员下的独立试标、规范冻结与金标放行}{fig:annotation}{\FigAnnotation}{已建立的标注工作流程。当前工作簿为空，流程不等于1500金标已完成；数量均为同一锁定主样本及其子集。}

家族ID、正文和公开日期随来源指纹保留，导入不能增加、删除或修改样本来凑配额。裁决记录绑定独立A/B原稿哈希，讨论后的统一答案不能重算成“原始一致性”。格式预检检查方向枚举、证据数量和正文逐字支持，人工审阅仍负责语义是否正确。只有完成裁决和明确放行的整批样本，才可以导出训练数据。

\section{QLoRA抽取器和正文表示}
\subsection{参数适配及独立质量评价}
Qwen3-8B采用NF4四比特量化和低秩适配，rank16、alpha32、dropout0.05，batch1、梯度累计16。LoRA和QLoRA分别提供低秩适配与量化微调方法\cite{r21,r22}，Qwen3提供中文语义基座\cite{r23}。完整序列上限2048 token，包含系统、用户和答案；正文预算从1024 token起测，不能把全部上限理解为可用正文长度。

截断或跳过超长记录会改变被评价语料，必须记录数量、类别、公开年份和失败原因。测试集只评价冻结模型，不能根据测试表现反复修改提示词或标注规则。抽取评价分别比较原生字段/规则、零样本和QLoRA，报告方向Macro-F1、原因短语评价、证据支持、严格JSON率及清洗后可解析率。解析合法率是接口质量，不是语义正确率。

原因和证据都是自由字符串，不能仅以字符串完全相等判定语义质量。方向可按类别统计；原因需检查是否概括了明确因由、是否混入结果；证据需检查覆盖和语义支持。模型自报置信度不作为正确率。独立测试的误判案例保留在私有记录中，用于解释下游分组结果，不把测试金标直接送入风险模型。

\subsection{正文池化、结构字段与缓存}
对每条去重公告只编码一次，正文向量为
\begin{equation}
 v_e=\frac{\sum_{q}m_q^{body}H_q^{last}}{\max(1,\sum_qm_q^{body})}.
\end{equation}
正文位置掩码排除系统提示、公司包装、padding和生成答案。若把生成JSON也池化，风险表示便混入模型自己生成的解释，无法清楚区分正文信息与生成反馈。图\ref{fig:pooling}说明池化范围。

\reportfigure{正文token与生成答案的池化范围分离}{fig:pooling}{\FigPooling}{拟议的冻结表示生成规则。正文预算和位置按分词器、模板版本验证，不展示真实公告内容。}

事件缓存同时保留预测三字段及原生数值。风险向量的数值结构包括方向编码、原因是否明确、证据可追溯状态和原生字段/缺失标记；自由原因短语与证据原文保留给解释和独立抽取评价，不把它们硬映射为预设行业类别。正文向量承担连续语义信息，风险投影从训练数据学习，不冻结随机矩阵。

\reportfigure{事件缓存键与日查询的版本关系}{fig:cache}{\FigCache}{缓存设计。每一输入或模型组成改变都会形成不同键，向量与预测字段使用同一可追溯版本。}

公司日事件聚合最多使用最近60市场日内的8条有效预告。对投影向量$\widetilde e_j$，内容分数经过数值截断后乘可用性、可追溯质量和年龄衰减，再归一化：
\begin{equation}
 b_j=\exp\{\operatorname{clip}(u^T\widetilde e_j,-20,20)\}m_jq_j e^{-r_eA_j},\quad
 e_i=\frac{\sum_{j\in\mathcal E_i(t)}b_j\widetilde e_j}{\max(\varepsilon,\sum_jb_j)}.
\end{equation}
全无有效事件时返回零表示并令$M_e=0$。质量$q_j$表达格式和来源可追溯状态，不等于人工语义正确概率；公司级$Q_e$按参与聚合的有效事件质量均值形成。若只使用二值质量，它主要起过滤作用，不能把$Q_e=1$解释为语义绝对正确。

诉讼5/20/60日计数、类型和最近年龄通过独立结构路径进入同信息模型。业绩预告缺失不应屏蔽诉讼事实；更新进展也不能回填首次公开日。由于Qwen预训练晚于回测期，本课题评价的是现代抽取器处理历史资料的回顾性效用。原生字段、仅预测结构和匿名文本控制有助于分析记忆风险，但无法证明预训练不存在污染。

\section{完整名单、金融主体与二部图}
\subsection{名单是更新单位，不是任意十条记录}
机构—公司关系依赖名单口径。原始前十大股东、流通股东、有限售股东、发行前股东和表决权名单描述的范围不同，不能把同一报告期的记录直接叠加。主图优先使用InfoTypeCd=1且OrgDisSHFlg=1的原始总股东候选；流通2/2和补充记录另行保存。SHTypeCd描述股东类别，不作为总股东/流通名单分组键。

完整性不能只看排名。新源中同排名的后续比较单位有54,743个，其中54,601个持股数原值相同，这提示大量合法并列的可能性。按相同排名只保留第一行，会删掉真实不同股东；任取十行，又可能混入补录或不同版本。拟结合原名单标题、报告期、公开日、来源、SHSer和ShrSer\_Type规则识别整组，而不是在未知组键上进行去重。

\reportfigure{从原始总股东候选到接受版本与隔离记录}{fig:certification}{\FigCertification}{名单与身份认证流程设计；严格结构候选只用于定位，不自动标为完整组。}

图\ref{fig:certification}把组认证和主体认证分开。新包缺少更新时间和披露版本ID，需要通过原版、更正及数据库历史覆写说明补证。至少100个披露组与100个主体的检查应覆盖跨年、并列排名、缺ID、产品更名和名义托管，而不是只挑容易匹配的案例。检查结果按问题类型记录，不能只给一个“抽查通过”总数。

\subsection{主体粒度与历史证券对应}
GDID有助于跨期识别，但不能先验地视为稳定受益所有人。一个标识出现多个原名称，可能涉及更名、产品版本或名义账户；一个名称对应多个ID，也可能是不同产品。基金产品与管理人并非同一个投资节点，名义托管更不能直接代表共同决策。个人股东、名义托管和不能可靠确认的主体从机构图中排除，记录排除原因和覆盖变化。

\reporttable{主体及名单更新中的处理取舍}{tab:entityrules}{
\begin{tabularx}{\linewidth}{p{3cm}YY}\toprule
资料情况&拟采用处理&需要保留的证据\\\midrule
ID与名称跨期一致&作为主体候选继续核验&原名称、类别、公开资料对应\\
同ID多名称&不直接合并为稳定受益人&更名、产品、名义账户规则\\
缺ID的非自然人&待类型与名称跨期确认&原报告与唯一对应范围\\
个人或名义托管&不作为共同机构关系&类型/托管识别与排除原因\\
新完整名单&按生效日替换旧整个列表&完整组、原版/更正版本键\\
股东退出Top10&移除当前披露图边&不把该行为当作真实清仓\\\bottomrule
\end{tabularx}}

公司历史证券对应采用独立有效区间，不能由DRET观察起止推断上市和退市。未可靠连接的辅助记录不进入当天图，目标公司日仍保留并表达$M_g=0$。这种设计避免“有图的公司较容易预测”被样本筛选误计为图模型提升，年度及共同覆盖子样本统计则用于解释适用范围。

\section{比例、边龄和目标去自身的消息}
\subsection{边权与机构聚合}
对接受的完整版本，边记录公司$i$、主体$j$、总股本持股比例$p_{ij}$和市场日年龄$a_{ij}$。主口径使用有限、非负的总股本比例，不把其他百分号字段一律当作相同量纲。定义
\begin{equation}
 w_{ij}=p_{ij}\exp\left[-\operatorname{softplus}(\lambda_{edge})a_{ij}/252\right].
\end{equation}
公司序列表示经过可学习变换$\phi$进入机构总和：$S_j=\sum_kw_{kj}\phi(s_k)$，总权$d_j=\sum_kw_{kj}$。向目标$i$返回时扣除自身：
\begin{equation}
 u_{j\setminus i}=\frac{S_j-w_{ij}\phi(s_i)}{d_j-w_{ij}},\qquad d_j-w_{ij}>0.
\end{equation}
分母不存在有效其他公司时，该机构不返回关系证据。图\ref{fig:bipartite}突出的是返回路径的扣除，而不是从整个图网络中删掉目标自身。

\reportfigure{二部图中目标去自身的加权消息路径}{fig:bipartite}{\FigBipartite}{仅为公司与机构变量的算法示意，不是实际股东个案，也不是因果传播图。}

进一步将机构消息按目标边权汇总：
\begin{equation}
 r_i=\frac{\sum_{j\in\mathcal N(i)}w_{ij}\mathbf1_{d_j-w_{ij}>0}\psi(u_{j\setminus i})}
 {\sum_{j\in\mathcal N(i)}w_{ij}\mathbf1_{d_j-w_{ij}>0}},\qquad
 g_i=\tanh\{W_g[s_i;r_i]+b_g\}M_{g,i}.
\end{equation}
目标自项$s_i$与跨公司消息联合投影是合法的图编码；需要去掉的是同一$s_i$先进入机构、再绕路返回的重复信息。若没有正权共享邻居，$M_g=0$。图随机化和同编码器对照可以检查$g_i$的增量是否主要来自参数容量或目标自项。

\subsection{小邻域和数值边界}
去自身后只有一个其他公司时，机构消息退化为该公司的变换表示。其边权在机构内归一化中抵消，不能据此声称持股比例在每种拓扑中都影响消息。存在两个或更多其他公司时，比例和年龄改变其相对贡献；外层不同机构之间还保留目标边权的作用。表\ref{tab:peercases}有助于区分这几种情况。

\reporttable{去自身后不同邻域规模的含义}{tab:peercases}{
\begin{tabularx}{\linewidth}{p{3cm}YY}\toprule
机构邻域&返回消息&对关系证据的解释\\\midrule
只持有目标公司&分母为0，不返回&不生成跨公司信息\\
仅一个其他公司&$u=\phi(s_k)$&机构内权重抵消，外层机构选择仍有作用\\
两个其他公司&二者的加权凸组合&比例和不同边龄改变相对贡献\\
较大邻域&分段加权聚合&需报告采样/截断及全邻居校验差异\\\bottomrule
\end{tabularx}}

实现中采用有效权重和数值阈值保护除法，记录极小分母、非有限值和权重下溢影响的边数。现有原型使用$10^{-8}$级阈值；正式数据预检应与更高精度聚合抽查，不能把数值保护阈值解释成经济意义上的最低持仓要求。重复公司—主体边应在来源处理阶段拒绝，而不是在图层里悄悄相加。

先计算各机构总和，再对每条目标边减去自身，可以使去自身消息量级保持为$O(Ed)$。若对每个目标都重新遍历全部机构邻居，高度机构会产生大量重复计算。研究将按日合并目标与邻居，使用当前权重编码每个公司一次，再进行分段聚合。训练采样上限可从15/10起测，正式验证与推断优先使用同一规则的分块全邻居计算。

\section{归一化边龄为何不足以限制整体修正}
对同一完整名单，目标公司多条边通常共享其披露年龄$A_i$。若外层边权为$p_{ij}e^{-rA_i}$，归一化后
\begin{equation}
 \frac{p_{ij}e^{-rA_i}}{\sum_h p_{ih}e^{-rA_i}}=\frac{p_{ij}}{\sum_h p_{ih}}.
\end{equation}
共同年龄因子抵消。即使整个目标名单变旧，归一化外层消息未必自动变小。邻居公司各自不同年龄仍会影响机构内相对权重，但这与限制整体辅助影响是两个问题。

这一现象是设置图外年龄包络的直接理由。图内权重处理“哪些邻居更值得参考”，图外包络处理“整体关系资料最多能改变主干多少”。若只把边龄写进归一化均值，可能遗漏目标名单绝对陈旧度。相反，若全部只在图外衰减，又不能区分不同邻居资料的新旧。两级设计需要分别消融，不能把任何含年龄变量的网络都称为同一种时点机制。

\section{有界残差融合及可检查性质}
\subsection{将修正幅度和内容分开}
主干分数由$s_i$与同信息模型共同采用的市场、行业和诉讼背景$c_i$计算，$z_{s,i}=w_s^T[s_i;c_i]+b_s$。仅收益GRU仍不接收$c_i$。公告和图风险头输出$z_e,z_g$，内容修正为$\delta_m=\tanh(z_m)$。对$m\in\{e,g\}$，设置
\begin{align}
 a_m&=M_mQ_m e^{-\operatorname{softplus}(\lambda_m)A_m}
 \sigma\{G_m([s;h_m;M_m;A_m;Q_m])\},\\
 z&=z_s+a_e\delta_e+a_g\delta_g,\qquad p=\sigma(z).
\end{align}
三路向量维度128，标量门输入维度为$128+128+3=259$。公告年龄取最近有效事件，按60市场日归一化；图年龄取目标公司当前接受完整版本，按252日归一化。事件注意力和事件残差共享事件衰减参数，图边及图残差分别使用参数，共三个年龄率。

$M_m$表示有效辅助信息是否存在。$Q_e$来自格式和证据可追溯检查，不使用未来标签或模型自报置信度；主图不可靠边在入图前排除，$Q_g$取1。门的内容项允许辅助修正为正或负，因此“有界”不意味着所有公告只会提高风险。

\subsection{幅度上界与同模型回退}
由于$0\le\sigma(G_m)\le1$且$|\tanh(z_m)|\le1$，有
\begin{equation}
 |a_m\delta_m|\le M_mQ_m e^{-\operatorname{softplus}(\lambda_m)A_m}.
\end{equation}
全无有效公告和关系时，$M_e=M_g=0$，从而$z=z_s$。这是同一模型的回退；它不等于另行训练、没有共同结构控制或辅助损失的收益GRU。图\ref{fig:envelope}以一个明确给定的衰减率说明允许范围的变化。

\reportfigure{辅助logit修正的解析年龄包络}{fig:envelope}{\FigEnvelope}{纯数学示意：$M=Q=1$且衰减率$r=1$，曲线为$\pm e^{-A}$，不是从训练资料拟合出的参数，也不是预测性能曲线。}

\reporttable{解析包络的数值及其解释}{tab:envelopetable}{
\begin{tabularx}{\linewidth}{p{2.3cm}p{3cm}Y}\toprule
归一化年龄 $A$&$e^{-A}$（示意$r=1$）&可以说明什么\\\midrule
0&1.0000&新鲜资料允许的最大修正\\
0.5&0.6065&幅度包络缩小，不确定实际方向\\
1&0.3679&不能推出概率必然下降\\
2&0.1353&不能当作已训练模型的真实贡献\\\bottomrule
\end{tabularx}}

门本身依赖内容和年龄，实际修正可能小于包络，也可能因内容变化而改变符号。因此，年龄上界单调不代表风险概率单调。由sigmoid导数上界$1/4$和均值定理，还可得
\begin{equation}
 |p-\sigma(z_s)|\le\tfrac14(|a_e\delta_e|+|a_g\delta_g|).
\end{equation}
该界描述同模型允许的概率偏移，不涉及泛化误差，也不保证限制后的排序更好。

\subsection{约束的代价和可检验问题}
有界修正会改变模型能表达的函数。如果某份旧资料仍有强信息，过度衰减可能损害预测；若事件质量标记只反映接口格式，还可能错误排除有用正文。这些代价不能由上述不等式排除。衰减率采用可学习参数，可以接近很小值，但能否学到合理权重仍取决于训练数据和优化。

研究将比较普通门控、年龄仅作特征和包络约束，检查自然年龄分组及人为增龄压力。压力实验固定资料内容和模型权重，只改变输入年龄，观察上界、分数偏移和排序变化；自然分组则保留真实资料状态，两者回答不同问题。普通门控也可以加入相似约束，所以方法区别最终仍需同编码器、同信息和合理预算下的结果支持。

\section{联合损失、训练掩码和优化}
风险任务主损失为训练目标公司日的未加权BCE。有效支路另加辅助损失：
\begin{equation}
 \mathcal L=\mathcal L_{main}+0.2\mathcal L_s+0.1\mathcal L_e+0.1\mathcal L_g.
\end{equation}
每个辅助BCE只在当前训练目标与该支路可用的交集上平均，没有交集时该项为0。仅收益模型只用主BCE；去辅助损失控制保留融合结构，移除额外支路监督。这样不会把邻居未来标签当成可共享训练目标。

\reportfigure{共享日批次的当前权重编码、消息计算与监督路径}{fig:training}{\FigTraining}{联合训练设计；反向传播可经过邻居历史表示，但不额外监督邻居未来标签。}

Qwen先完成抽取评价再冻结；正文原始向量可缓存，事件投影、GRU、图消息与门在风险任务中共同学习。冻结随机GRU或随机图投影会让输入表示与任务不匹配，不能以“节省训练”代替正式联合学习。训练采用AdamW，学习率候选$3\times10^{-4}$和$10^{-3}$，weight decay为$10^{-4}$，最多20轮，验证AP连续3轮不改善早停，梯度范数上限1。

概率校准只在规定验证集拟合温度或截距，检查Brier及校准情况\cite{r30}。在校准与稳定性评价之前，输出宜称风险分数。所有配置、预处理参数、缓存键和随机种子与运行记录一起保存，便于追踪模型选取及恢复后的结果。

\section{时间外划分与同信息强对照}
\subsection{扩展窗口和结果区间剔除}
测试折依次为2021、2022、2023和2024年截至7月24日，验证使用测试前一年，训练窗口逐年扩展。图\ref{fig:folds}只按年格呈现结构，实际索引按市场日和结果结束日执行边界检查。训练标签结果结束日必须严格早于验证首个市场日，验证标签结果结束日必须严格早于测试首个市场日。

\reportfigure{四段扩展训练、独立验证与时间外测试}{fig:folds}{\FigFolds}{已生成的四折共享索引设计。年格不表示精确日长度，2024只覆盖至7月24日，purge依据结果结束日。}

公司日索引只有一份，模型通过信息权限读取不同输入。训练序列可以向前读取验证日前的历史，不能向后读取结果区间；关系邻居使用预测日已知的历史表示。缺失公告或关系时，保留目标标签和序列，用掩码表达辅助不可用。另报共同覆盖子样本有助于理解关系稀疏程度，但主要比较仍使用共同全量目标。

时间序列交叉验证的有效性依赖具体条件\cite{r33}，不能由某些条件下可用随机划分推出股票公司日也可以打散。市场共同冲击、重复公告家族和五日标签重叠同时存在，独立年份与结果区间剔除更贴近这里的预测问题。预处理、文本降维、校准与选模各自限定拟合集合，避免模型在测试分布上调整后再宣称时间外表现。

\subsection{比较组和输入预算}
收益-only组包括下行波动/传统尾部统计Logit、强树和GRU。历史特征可含滚动残差偏度、上下行波动比、波动和回撤，全部由起点已知资料计算。LightGBM、XGBoost、CatBoost是成熟非线性表格工具\cite{r25,r26,r27}，正式比较应给予合理验证预算，而不是仅使用默认参数的弱线性模型。

同信息组取得原生及预测事件、正文聚合、诉讼、关系统计、边龄和邻居历史量。强树取得同一正文向量的训练期PCA64表示，不能只给它公告数量，却让主模型独享语义信息。同编码器组则保持同GRU、事件缓存和二部消息，主要改变拼接、固定融合、普通门控、MSGCA式交叉注意力适配和本方法。

\reporttable{比较模型的信息权限及适配说明}{tab:baselines}{
\begin{tabularx}{\linewidth}{p{3.05cm}YY}\toprule
模型组&信息权限&比较目的\\\midrule
收益统计Logit&收益、三因子及训练期历史统计&量化简单风险特征能力\\
收益强树/GRU&同收益主信息&检验序列表示与非线性表格差别\\
同信息强树&同正文、事件、诉讼和图统计&排除信息剥夺造成的虚假增量\\
同编码器融合&同三路表示和结构控制&检验融合约束，而非换编码器\\
图时序适配&同标签、日期与允许关系&研究直接图方向在本任务下的表现\\\bottomrule
\end{tabularx}}

Graph WaveNet及多关系风险方法按同一五日标签适配。直接近邻完整方法或代码尚未取得时，只能标为启发式适配，不能称忠实复现。每折主对照从全部同信息模型中按验证种子平均AP选择，排除本研究模型M；测试不参与选择。树模型每信息组每树最多8个验证配置，神经模型采用共同候选和停止规则，保留全部选参记录。

\section{风险排序、概率质量和配对统计}
\subsection{指标的分工}
AP用于衡量正类风险排序。Davis和Goadrich以及Saito和Rehmsmeier讨论了PR与ROC在不平衡评价中的关系\cite{r28,r29}。本数据P10正类占比约10.94\%，仅凭多数类就能得到较高分类准确率，因而准确率无法充分说明有限审核名单的价值。AP也受正类率影响，需逐折报告实际正类率、AP和差值。

每日前10\%取$k_t=\lceil0.1N_t\rceil$，分数同分时按证券代码固定排序。报告每日精确率及有正例日期的召回均值，而不是把全年公司日混排后选前10\%。F1阈值和温度校准只在验证集拟合；ROC-AUC和Brier作为辅助评价，与主排序增量分开解释。模型概率在未经验证校准以前只称风险分数。

\subsection{种子与日期不确定性分别计算}
双方每折固定种子42、2024、7、123、3407。令$AP_{f,s}^{M}$和$AP_{f,s}^{B}$为模型与该折验证选出主对照的指标，主差值为
\begin{equation}
 \Delta AP=\frac14\sum_{f=1}^{4}\left(\frac15\sum_{s=1}^{5}AP_{f,s}^{M}-\frac15\sum_{s=1}^{5}AP_{f,s}^{B}\right).
\end{equation}
先计算各模型各个种子AP，再取均值。AP对排序是非线性的，所以$AP(\frac15\sum_sp_s)$通常不等于$\frac15\sum_sAP(p_s)$；概率集成可以另报，但不是同一个估计对象。

\reportfigure{同公司日配对、种子均值和日期块区间的计算路径}{fig:evaluation}{\FigEvaluation}{统计方案，不包含模拟预测分数或区间。日期块在折内抽取，再进行四折等权汇总。}

K\"unsch和Politis、Romano的工作为依赖序列重采样提供了基础\cite{r31,r32}。本课题采用整市场交易日的配对移动块，20日为主块长，5/60日为敏感性，1000次。一次抽取中，同一天的全部公司及两模型预测一起进入样本，保留横截面共同冲击。这里不是平稳bootstrap的随机长度算法原样复现，具体抽样规则须与实现记录一致。

训练种子标准差刻画优化不确定性，日期块区间刻画样本时序不确定性，两者分别报告。四折不按样本量加权，以免公司日更多的年份支配结论。主对照选择发生在验证期，区间计算发生在冻结测试预测后；多次试探和事后选择会影响解释，White关于数据窥探的讨论对此有直接警示\cite{r34}。

\section{十一项消融与机制压力}
消融区分三个问题：信息是否存在，限制是否有效，结果是否来自不可靠来源或关系巧合。表\ref{tab:ablations}中的11项保持共同样本和训练协议；P5作为稳健性标签，普通门控和拼接属于架构对照，不混入这11项计数。

\begin{longtable}{p{.65cm}p{3.1cm}p{4.5cm}p{5.3cm}}
\caption{固定消融、改变内容和所检验问题}\label{tab:ablations}\\
\toprule 序号&实验&改变内容&检验问题\\\midrule\endfirsthead
\toprule 序号&实验&改变内容&检验问题\\\midrule\endhead
1&去公告&关闭预告分支&公告信息是否提供额外效用\\
2&去图&关闭关系分支&关系系统整体增量\\
3&去全部年龄&关闭事件注意力、边消息及修正年龄&年龄处理整体作用\\
4&去质量&移除质量过滤/加权&接口质量策略是否有效\\
5&年龄仅作特征&保留年龄输入，无幅度包络&信息变量与幅度限制差别\\
6&去辅助损失&只用主任务BCE&分支监督与融合约束的作用\\
7&仅去边级衰减&保留图外包络&相对邻居权重作用\\
8&度保持随机图&随机重连关系&真实关系与容量/度数效应\\
9&仅原生公告&不使用Qwen语义&语言适配的必要性\\
10&仅Qwen结构&移除正文隐藏表示&结构抽取与连续语义差别\\
11&匿名文本&匿名公司/证券包装&预训练记忆风险的敏感性\\\bottomrule
\end{longtable}

去质量控制可能同时改变有效事件集合，不能把差值解释为纯粹标量门效应。随机图应保持度数约束、预测日可见性和机构独享规则，不把重连引入的非法未来边当作关系随机性。匿名化仅减轻记忆线索，无法证明模型未记住事件内容。每项实验需将变更的数据产品、参数和输入权限列入记录。

自然分组按模型输入中可获得的年龄、质量和共享邻居规模划分，观察新旧事件、低覆盖公司和不同市场年份。人为增龄保持内容与权重固定，检查允许上界与实际分数变化。这种压力检验用于解释机制，不作为新测试样本或额外金融收益结论。若总体增量仅由某一年或高覆盖群体驱动，需明确适用条件。

\section{日装配、存储与训练恢复设计}
\subsection{共享数据而不是重复宽表}
把6,116,554个标签全部展开成60日、五字段float32窗口，仅数值本身就约为
\begin{equation}
 6{,}116{,}554\times60\times5\times4\approx7.34\times10^9\text{ 字节}.
\end{equation}
它不包含掩码、索引、正文向量、图或运行时副本，已接近当前主机内存尺度。研究拟保存原始序列与位置索引，通过市场日查找窗口；事件缓存全局共享，不按公司日重复保存。上述数值是存储理论下限，不是已经测得的程序峰值。

日装配器查询目标和邻居的并集，统一标准化后返回序列、事件、边及监督掩码。当天模型计算一次公司表示，再用稀疏索引聚合机构。训练目标集合和计算节点集合必须分开；某公司作为邻居出现，不自动获得未来标签监督。跨日期合并若改变图、年龄或控制变量，则不能简化为一个静态全样本图。

\reporttable{数据产品及版本失效条件}{tab:storage}{
\begin{tabularx}{\linewidth}{p{3.1cm}YY}\toprule
产品&保存内容&何时失效或需重建\\\midrule
标签及四折索引&日期、证券、结果结束日、标签、split&收益/因子、标签算法或范围变化\\
原始序列索引&市场位置、字段及掩码&字段顺序、日历、标准化版本改变\\
正文缓存&原文哈希、模型/模板、正文位置、预测字段&任一缓存键组成变化\\
完整图版本流&公司、主体、比例、公开版本、生效日&名单/主体/证券对应规则修订\\
批次manifest&输入版本、维度、允许信息、目标集合&任何产品或装配合同变化\\
冻结预测&模型、索引、种子、日期与输出&模型/输入改变，不覆盖旧结果\\\bottomrule
\end{tabularx}}

\subsection{复杂性与计算瓶颈}
设$N$为单日去重公司数、$J$为机构数、$E$为边数，序列长度$T=60$，实际GRU输入通道$F^*=10$，隐藏维$h=64$，投影$d=128$。序列计算约为$O(TN(F^*h+h^2))$；二部变换及聚合约为$O((N+J)d^2+Ed)$；事件聚合受$K\le8$限制。编码当前公司一次后再共享给各机构，避免重复GRU成本。

大模型正文编码成本发生在事件级缓存阶段，不能隐藏在风险网络每步耗时之外。资源日志分别记录抽取、缓存、装配、风险训练和评价I/O。高度机构采用分块全邻居聚合；采样策略先在小批次比较输出偏差和截断率，不能只报告节省内存的结果。实际小时数、显存和费用在GPU试跑后测量，不用合成CPU步时外推。

\subsection{恢复运行与异常处理}
检查点绑定模型配置、优化器、随机种子、输入版本和训练位置。恢复后不得重新选择验证表现更好的种子或合并不同缓存；预测按共同索引检查重复、漏行和年份边界。记录可重复命令和环境，特别区分CPU与GPU依赖、量化库及模型revision。

\reporttable{常见失败、影响与定位方法}{tab:failure}{
\begin{tabularx}{\linewidth}{p{3.15cm}YY}\toprule
失败情况&可能影响&定位或处理\\\midrule
代码命名空间混用&公告/边连到错误证券&源字段追溯，唯一历史区间确认\\
名单口径混合&重复边及假共享关系&类型分离、原版比对，不按排名删行\\
历史覆写不可追溯&日期正确但内容有未来信息&版本政策与原公告，未决资料隔离\\
正文范围或截断变化&缓存包含答案或与模型不一致&token位置检查及缓存键失效\\
真实批次内无监督目标&错误训练或空损失&立即失败，不用邻居标签补目标\\
恢复后漏预测/重复索引&配对评价样本改变&全量键检查，按冻结索引重读\\
年龄约束过强&有效旧信息被削弱&普通门控、自然分组和增龄压力比较\\\bottomrule
\end{tabularx}}

异常处理应使问题可见。原版匹配失败、缺少ID、解析失败或极小图分母都需要独立计数，不能用一个“缺模态”标签掩盖全部原因。已完成输入检查并不意味着标签或主体永远正确，后续修订必须保持旧数据版本和结果，重新确认受影响实验。

\ReportTypesettingCheck\FloatBarrier
\chapter{研究工作进度安排}
\section{排期依据和主要依赖}
研究排期围绕资料、金标和真实输入组织，而不是按模型名称依次运行。收益标签和共享四折已经提供底座，当前最先影响后续训练的是名单/主体/历史版本证据与人工预告金标。日装配器可在CPU先建立接口和边界测试，但真实Qwen缓存需要抽取模型评价和冻结后才生成。这些依赖决定了哪些任务能够并行，哪些只能等待前置条件。

\reportfigure{来源证据、人工标注、CPU接口与正式实验的依赖}{fig:dependencies}{\FigDependencies}{研究实施设计；箭头表示产物依赖而非已完成比例。来源核验和人工标注可与CPU接口并行，正式训练在前置条件满足后启动。}

图\ref{fig:dependencies}中的并行部分具有现实意义。等待数据库回复时，可进行真实窗口、事件版本流、缺失掩码及索引检查；等待标注裁决时，可完成普通门控和强树的信息接口。另一方面，随机文本向量或未认证边不能用于正式比较，不能为了赶排期先跑一轮再把这些替代输入写成实证。

\section{目标阶段和可检查产物}
表\ref{tab:schedule}使用目标时间，不代替学院校历。学校两年制专业学位硕士开题应在第三学期第1周（含）前完成，用户入学年级和正式考核安排尚未提供，因此本稿的编制日期不能证明该手续按时完成。研究工作安排由导师结合个人培养计划复核。

\reporttable{阶段工作与产物安排}{tab:schedule}{
\begin{tabularx}{\linewidth}{p{2.45cm}YY}\toprule
目标阶段&主要工作&提交或核查产物\\\midrule
2026年11月底&补充开题正文、图表与文献；定位来源证据&导师可审稿、文献阅读记录、原报告任务表\\
12月上旬&名单/主体/证券核验，200/42独立试标；CPU真实接口&原版问题分类、裁决及规范版本、窗口/时点测试\\
12月中旬&完成剩余1500/300标注；共同信息对照和恢复调度&人工放行数据、SFT整批预检、日装配与配置\\
前置条件满足后&GPU长度/显存/吞吐试跑，抽取器独立评价和缓存&冻结模型与缓存键、资源日志、真实输入核查\\
目标12月下旬至1月&四折双方五种子主比较、固定消融及P5&验证选择台账、全部种子结果、配对统计\\
目标2027年2月&整合论文、复核引用数字、导师修改与答辩&完整论文及非敏感复现材料；正式截止依学院\\\bottomrule
\end{tabularx}}

人工标注约1800次判定加裁决，工时应从200条试标记录估计。300双标是主样本子集，不额外增加300篇公告。正式全矩阵还包含对照验证配置、各种子重训、消融和P5，不能只按“一个模型一轮”计算时间。GPU小时、I/O等待和中断恢复分别计入，再预留运行余量。

\section{排期偏离时怎样调整}
若历史版本无法取得，首先列清无法认证的组、所占年份和覆盖变化，明确可接受与隔离范围；不把最新修正版自动当作历史原版。若金标分歧较多，先修订训练试标规范并补裁决，而不是提前把不一致答案送入SFT。若真实吞吐低于预期，先检查窗口重复、邻居重复编码、缓存I/O和恢复粒度，再调整计算批次及租赁时段。

这些调整改变的是工程执行和资源安排，不应改变共同测试样本、主指标或测试选模规则。方案保留完整三路编码、关系消息和有界融合，实验结果可以否定部分假设；不能在看到失败结果后删掉弱年份、种子或对照，再把剩余结果称作预先确定的实验。

阶段成果还应与培养手续区分。导师内容审核、政治立场和价值导向审核、开题口头考核及系统上传需要本人和培养单位实际办理。文稿、程序和实验产物为这些环节提供材料，不自动产生审核意见、签名或考核结果。
\ReportTypesettingCheck\FloatBarrier
\chapter{预期研究成果}
\section{方法成果与证据对象}
本课题拟形成一个由标签、输入、模型和评价相互约束的短期尾部预测方法。标签明确公司日及五日结果，输入明确实际公开日期和主体，模型给出收益主干与辅助修正，评价记录年度和随机种子。这些对象共同决定一份风险排序的含义，单独保存网络权重不足以复现研究。

方法成果包括三字段抽取、按公开时点更新的完整关系版本、目标去自身消息及有界残差融合。抽取成果既包含模型也包含独立语义测试；图成果既包含边也包含来源和身份排除记录；融合成果既包含公式和代码，也包含性质、压力和同信息比较。它们分别回答语义表达、关系可信性和修正限制三个问题。

\reporttable{预期成果、使用对象及对应证据}{tab:outputs}{
\begin{tabularx}{\linewidth}{p{3.2cm}YY}\toprule
成果&用途&证据对象\\\midrule
条件尾部标签与四折索引&共同定义模型目标&算法版本、日期边界、分布与独立参照\\
公告规范和抽取模型&从正文获得有限结构&独立原稿、裁决、语义指标与解析失败\\
完整持股版本与主体表&重建预测日关系&原报告、类型/身份、覆盖与年龄\\
有界风险模型&学习受限制的辅助修正&性质推导、梯度检查、同编码器对照\\
正式比较与消融&检验信息和机制增量&全部种子、年度结果、配对区间和失败分析\\
可追溯离线原型&支持有限审核名单研究&输入/模型版本、主干及两路修正、异常记录\\\bottomrule
\end{tabularx}}

\section{风险排序与追溯原型}
图\ref{fig:output}展示拟议输出关系。公司日分数与当日排序名单分开，名单容量按当日可预测公司数计算。结果同时保存主干、公告修正和图修正，让研究人员能够追查高分来自价格状态、事件还是关系。解释这些项时，不把权重和修正方向认作公告的因果作用。

\reportfigure{风险排序、分支修正和版本追溯的输出设计}{fig:output}{\FigOutput}{拟议原型，不含实际公司预测或模拟风险数值。输出为离线研究与审核辅助，不是交易建议。}

原型需能回答一组简单但关键的问题：该行使用哪个模型与索引版本？输入中最晚可用资料是什么？公告是否因解析失败被排除？名单是否存在共享邻居？分数是否经过验证集校准？这些信息不改变AP公式，却决定结果能否被复核。未映射或无辅助资料的公司不应在导出时丢失，而应保留状态并回到相同模型主干。

公开交付以论文、方法代码、配置、数据合同、聚合统计和复现说明为主。授权原始记录、公告正文、逐样本预测、适配器和风险权重按许可留在私有环境。源码包和报告图表不得夹入旧论文正文、个人签署声明或原始数据。

\section{成果判定与论文组织}
论文拟按问题与相关研究、数据时点、方法、实验及总结组织，开题阶段使用学校七栏目描述研究计划。最终主表由冻结预测生成，包含模型、信息权限、测试年、样本量、正类率、AP、差值、区间、每日名单指标及资源消耗。原因/证据抽取与下游排序分表报告，不能用一个较好的抽取分数代替金融预测结果。

如果增量不能稳定成立，研究仍需给出完整实验及失效条件，分析覆盖、抽取误差、年龄分组与市场漂移。成果质量由方法差异、证据完整性和学院评阅决定，不由预先承诺的固定AP增幅决定。是否形成投稿或专利申请，以实际结果及导师判断为准，不将发表等级当作已经获得的成果。
\ReportTypesettingCheck\FloatBarrier
\chapter{本课题创新之处}
\section{从年龄变量到辅助修正幅度的约束}
收益、文本、图和普通门控都已有研究基础。本课题拟检验的第一项差异，是在风险logit层明确区分主干和辅助修正，并把可用性、质量和年龄放入修正幅度的乘积包络。年龄只是普通输入时，预测头可以学习任意非线性关系，却未必满足资料陈旧时的最大作用限制；有界结构则将这一性质直接写进模型。

这种差异同时带来约束和代价。它使全缺失状态有明确回退，并提供随年龄缩小的允许偏移，但也会缩小函数类。有效旧信息可能因衰减而受到限制，解析质量标记也可能漏掉仍有信息的正文。研究贡献因此不是一个“必然提高”的不等式，而是检验结构限制在目标任务中的净效用。

\reporttable{年龄普通特征与有界修正的可检查差别}{tab:mechanismcomparison}{
\begin{tabularx}{\linewidth}{p{3.3cm}YY}\toprule
比较维度&年龄作为普通特征&年龄有界修正\\\midrule
资料陈旧度&由预测头任意利用&进入可学习非负衰减包络\\
最大辅助偏移&一般无显式年龄界&$M Q e^{-rA}$为分支上界\\
全缺失状态&依具体偏置/归一化实现&修正为0，回同模型主干\\
新鲜资料的利用&表达自由度较大&保留内容门和双向tanh修正\\
待测代价&噪声可能被放大&可能损害有效旧信息\\\bottomrule
\end{tabularx}}

表\ref{tab:mechanismcomparison}不是对所有门控方法的普遍定理。普通门控也可以被加入同样约束，因此真正需要比较的是具体实现、共同信息及编码器预算。提出的方法差异必须在同编码器控制、年龄仅作特征、去全部年龄和去质量实验中形成证据，而不是仅凭结构图认定创新。

\section{相对邻居权重与绝对陈旧度的两级处理}
归一化关系消息存在一个容易忽略的性质：目标同一名单的共同边龄因子可以抵消。仅写边龄权重，可能只改变不同邻居的相对贡献，却没有限制整份目标名单变旧时的整体影响。本课题将边级衰减和图外修正包络分开，前者处理邻居相对新旧，后者处理目标关系资料的绝对年龄。

去自身聚合进一步避免目标表示通过机构均值返回，独享机构不提供跨公司证据。目标自身仍可在图投影中与合法消息联合编码；关系随机化和同编码器预算用于检查增量是否来自容量，而非真实关系。此项属于面向本任务的消息设计和公平实现，不把GraphSAGE或去自身思想单独宣称为首创。

单一其他邻居、全部目标边同龄和高度机构是三个有价值的检验边界。前者使机构内权重抵消，第二种情况使外层年龄因子抵消，第三种情况产生资源与采样问题。这些边界使公式、数值检查和实证设计相互对应，也给“年龄究竟在哪一层起作用”提供具体解释。

\section{可追溯的工程合同与信息公平验证}
不同资料只要在名称和日期上能够拼接，并不意味着其经济对象和历史可用内容一致。本课题拟把来源、版本、主体、缓存和日索引贯穿到预测导出中，模型和强对照共享同一公司日。预告抽取与风险预测分开评价，正文信息同样给予强树，避免把输入差别当成融合增量。

这一工程合同的实践价值在于，结果不再只是一列概率。它能够说明某个分数使用了哪份公告、哪份名单、何种模型、哪些预处理参数，以及哪些资料因证据不足未进入。来源核验、全文语义评价和模型统计各自形成独立记录，在异常或版本修订时可以判断需要重做哪些步骤。

\section{创新主张需要满足的证据条件}
近期图时序和MSGCA已经覆盖与本方案接近的骨架，相关专利CN117808586A的公开线索也涉及风险目标和多关系融合。完整权利要求及部分近邻全文尚待核验，不能由摘要未出现某个词就推断方法不同。后续比较需逐项核对标签频率、资料可用性、关系重构、融合形式和评价协议。

若有界结构只相对收益-only模型提高，而不能超过相同信息普通门控或强树，则主要证据是新增信息而非融合创新。若增量只由单一年份或高覆盖群体驱动，应限定适用场景；若区间跨零，应报告未获得稳定增量。只有方法差异、性质、同信息比较和机制分析一致，才支持更强的研究贡献表述。投稿竞争力还涉及近邻充分性和真实结果，不由开题阶段预先保证。
\ReportTypesettingCheck\FloatBarrier
\chapter{研究基础}
\section{已取得的资料与数据质量分析}
\subsection{收益、标签与共同索引}
日收益9,091,641行，观察范围2015年1月5日至2024年7月31日，全市场日期2329个。有效标签6,116,554行，范围2017年2月6日至2024年7月24日，P10正类669,391、P5正类358,973。标签规模与收益行数不同，既涉及历史回归及完成窗口预热，也涉及未来五日结果完整性，不能将其差值全部归因于单一缺失字段。

\reporttable{已验数据底座及对应统计口径}{tab:foundationdata}{
\begin{tabularx}{\linewidth}{p{3.6cm}p{3.2cm}Y}\toprule
对象&规模或日期&统计口径\\\midrule
日收益&9,091,641行&2015-01-05至2024-07-31\\
全市场日历&2329个日期&不是每只股票的有效收益日\\
有效公司日标签&6,116,554行&2017-02-06至2024-07-24\\
P10正类&669,391&同一有效标签面板，约10.94\%\\
P5正类&358,973&同一面板稳健性标签，不另加公司日\\
业绩预告&162,168条&原事件总量，模型输入另按公开日期查询\\
诉讼仲裁&21,219条&行为描述可用，最新进展不能回填首次日\\\bottomrule
\end{tabularx}}

全量另目录重建标签与旧表解压内容一致。独立NumPy参照抽查20家公司19,476标签，残差、阈值和P10/P5无数值或标签不一致。该参照是抽样，不是第二套算法独立全市场重算。共享四折索引已全量读回，标签、ID、结果结束日及split边界一致，提供了正式模型共享目标的底座。

\subsection{新股东补导的公开时间}
最新原包3,607,689行、76字段，前三工作簿各1,048,575数据行，第四461,964行。外层及工作簿内层CRC通过，输入前后SHA不变；独立计数与原值CSV全量来源读回一致，另用独立XML解析抽查16源行的76字段，原值一致。旧4,036,775行及原包仍保留，新旧源历史逐行对齐尚未完成，不以行数减少宣称去重成功。

报告期范围为2015-01-05至2024-12-31，公开日期却延至2026-07-16。图\ref{fig:yearcounts}按源记录的公开年份统计，2024年包含该年7月24日以后的记录，2025/2026也属于原包；它们不能都进入固定风险样本。研究信息观察窗口内公开记录为3,297,707条，窗口外为309,982条。

\reportfigure{新股东原包记录的公开年份分布}{fig:yearcounts}{\FigYearCounts}{新包全量审计聚合。各年计数合计3,607,689；这是源公开年份而非训练样本年份，2024全年和2025/2026记录需另按2024-07-24截止过滤。}

\begin{longtable}{p{2.4cm}r p{2.4cm}r}
\caption{公开年份图的数值表}\label{tab:yearcounts}\\
\toprule 公开年份&记录数&公开年份&记录数\\\midrule\endfirsthead
\toprule 公开年份&记录数&公开年份&记录数\\\midrule\endhead
2015&200,294&2021&392,168\\
2016&284,746&2022&408,995\\
2017&314,438&2023&430,988\\
2018&333,991&2024&427,881\\
2019&345,614&2025&101,705\\
2020&366,529&2026&340\\\bottomrule
\end{longtable}

2025/2026的晚公开记录并不自动说明原包错误，它们可能描述旧报告期的补充或更正，但也暴露出“按报告期筛选历史表”不足以重建当时信息集的问题。35,192条报告截止日晚于公开日也仅列为诊断，需核业务含义，不能直接认定为错行或提前生效。

\subsection{字段缺失与名单结构}
图\ref{fig:missing}使用同一个全量行数分母，展示A股代码、股东ID、流通股比例和来源缺失。股东ID大量缺失包含自然人，流通股比例缺失也与名单口径有关，因此两类缺失不能直接解释为机构图损坏。主图使用总股本比例，PctTotShr在新源中全量数值有效且位于0—1；本次范围检查不等于独立全量重算股份分母。

\reportfigure{新股东原包的字段缺失比例与数量}{fig:missing}{\FigMissing}{分母均为3,607,689条源记录；计数由全量审计得到。图不表示机构身份错误率或公司日模态覆盖率。}

\reporttable{字段与身份诊断的不同分母}{tab:missingdata}{
\begin{tabularx}{\linewidth}{p{3.6cm}p{3.25cm}Y}\toprule
诊断&数量&分母和解释\\\midrule
A股代码缺失&9,318条&全包记录；不自动回填历史证券\\
GDID缺失&1,732,902条&全包记录，包含自然人等\\
企业属性行缺GDID&1,289条&企业属性1,271,367条，不等于金融机构总数\\
流通股比例缺失&2,021,545条&全包记录，不能据非空认证流通名单\\
信息来源缺失&33条&全包记录，需定位原始来源\\
同GDID多原名称&9,018个ID&观察多对多，不直接合并受益主体\\
同原名称多GDID&69个名称&不能以名字直接创建稳定节点\\\bottomrule
\end{tabularx}}

原始总股东1/1候选全期180,652组，研究信息窗口内165,192组。严格十行及十个唯一排名条件下，全期142,838组、窗口内129,500组满足结构筛选。这个条件没有包括可能合法的并列排名，所以不能当作全部可用完整名单数量。33,597组存在重复排名，后续比较54,743个单位中54,601个持股数原值相同，原报告和序号规则因而比单纯去重更重要。

新包没有UpDt、nobs、ID或DisclosureID，下载的词典HTML仅为官方跳转入口；主要名单枚举已有用户提供的官方页面取值，ShRank并列、ShrSer\_Type分组、GDID稳定范围及历史覆写规则仍待正文说明。结构、身份、公开版本和证券对应必须分开核验，当前\texttt{main\_graph\_ready=false}。

\section{程序和研究积累}
已有标签、四折、来源映射、中文标注、裁决和SFT整批预检工具。来源篡改、成员错误、证据不在正文、未完成裁决等情况会被拒绝。模型原型已经测试缺失回退、修正上界、梯度、去自身和独享机构等性质。77项Python合成功能回归包含数据、模型、标注、映射及学校文稿检查，另有PPT源检查，不能将其总数当作金融效果实验数。

旧11特征先导Logit平均AP为0.221784，LightGBM为0.236579，差约1.48个百分点。这属于既有表格比较，不是本研究M的提升；旧先导的独立验证、时序区间和每日名单口径还需按新协议重做。本报告没有正式Qwen、真实持股图或完整风险模型的实证分数。

\section{实验条件、缺口和取得途径}
当前环境约4CPU、8.8GiB内存，无CUDA，可进行源审计、接口和小样本测试。收益、三因子、预告、诉讼与股东授权资料已有，不把未核价格、市值、换手或财报时点列为必要输入。真实60市场日取窗、事件版本流、日装配器、若干同信息对照和全矩阵调度仍待落实。

人工金标目前为0。中文试标工作簿与私有包已经生成，但还未传到本人本机完成实际标注。下一步由A200/B独立42形成初始反馈，按规范完成其余标注和裁决。原版名单、更正和主体通过新版10组定位，再扩展到至少100组/100主体；受限词典、历史版本和证券有效日期由可授权渠道补证，不要求手工清洗全市场原表。

GPU尚未租用，Qwen与正式风险模型未训练。拟在来源、金标和CPU接口检查满足后，使用实际设备测完整序列显存、吞吐、恢复和I/O。24GiB是起始参考，不是已经保证容纳的条件。没有已经落实的国家重点实验室或企业实验室合作，不虚构实践经历或外部支持。

\section{经费与复现安排}
经费包括资料补证、人工标注、GPU抽取/缓存、风险全矩阵及备份。现有授权数据和CPU优先复用；增量费用归属、GPU报价和总小时待核。预算以实测小时$H$、时费$c$及必要余量计算，人工工时由试标记录估计。学生计划自行租算力，但未提供金额、批复或付款证明，具体经费来源及落实状态仍待确认。

原始资料只读，新加工和预测保留在授权私有目录。公开报告和源码只含方法、图示、聚合事实和自写文稿；原公告、股东原行、样本预测与权重不公开。当前学制为两年专业学位硕士，学院、类别/领域、导师、学号、入学年级和正式开题日期未给出，封面保持待填写。本人实际文献研读、导师审核、实践环节和系统手续仍须按个人培养方案办理，不能由程序或文稿代填完成状态。
\ReportTypesettingCheck\FloatBarrier
\clearpage
\edef\ReportBodyPages{\number\numexpr\value{page}-1\relax}
\typeout{REPORT_BODY_PAGES=\ReportBodyPages}
\backmatter
\input{references.bbl}
\end{document}
